\chapter{Right-Perfect Morita Theory and Canonical Morita Duality}
\label{chap:right-perfect-morita-duality}

\section{Purpose and standing conventions}
\label{sec:right-perfect-morita-purpose}

Fix
\[
    A\in\CAlg(\V)
\]
and write
\[
    \mathcal W
    :=
    \Mod_A(\V).
\]
This chapter develops the noncommutative Morita theory of associative
algebra objects in
\[
    \mathcal W.
\]
Its objects are associative $A$-algebras, its $1$-morphisms are
bimodules, and its $2$-morphisms are bimodule maps. Horizontal
composition is relative tensor product.

The passage to Morita theory does not alter the algebraic foundations of
the preceding chapters. It raises the categorical level by one.
Dualizable objects are replaced by adjointable $1$-morphisms, ordinary
duals are replaced by Morita adjoints, and the canonical rigid duality is
replaced by the composite of right-adjoint formation and action reversal.
All algebraic constructions continue to take place in the presentably
stable symmetric monoidal category
\[
    \mathcal W
    =
    \Mod_A(\V).
\]

The principal task is to isolate the bimodules admitting Morita right
adjoints and to construct the coherent duality
\[
    B\longmapsto B^{\op},
    \qquad
    M\longmapsto(M^R)^{\op}.
\]
The two operations in the second formula have different variances. Right
adjunction reverses the direction of a Morita $1$-morphism, and action
reversal reverses it once more. Their composite therefore preserves the
direction of $1$-morphisms while reversing $2$-morphisms.

The chapter separates the following layers:
\begin{enumerate}[label=(\roman*)]
    \item associative algebras, bimodules, and the higher Morita
    $(\infty,2)$-category;

    \item one-sided perfectness and its equivalence with Morita
    adjointability;

    \item opposites, elementary mates, and the coherent mate calculus;

    \item specified adjunctions and the invariant mate--reversal
    construction;

    \item canonical Morita duality and its compatibility with base
    change.
\end{enumerate}

The conventions of
\Ref{chap:relative-ag-monoidal-syntax},
\Ref{chap:perfect-rigid-enrichment},
and
\Ref{chap:hermitian-poincare}
remain in force.

\begin{convention}[Higher coherence in the Morita layer]
\label{conv:rigid-morita-higher-coherence}
Adjunction, mate, action-reversal, pseudofunctoriality, biduality, weak
involution, and base-change data are always understood together with all
required coherent higher homotopies.

Whenever such data are suppressed from the notation, their existence is
supplied by a coherent adjunction or mate theorem, by a universal property
with contractible choice space, by transport along an equivalence, or by
an explicitly cited coherence theorem.
\end{convention}

\begin{convention}[Invariant $(\infty,2)$-categorical language]
\label{conv:infinity-two-variance}
We work invariantly with $(\infty,2)$-categories and functors between
them. No particular presentation of the homotopy theory of
$(\infty,2)$-categories is fixed.

Whenever a cited construction or coherence theorem is proved using a
particular model, only its invariant conclusion is used in this chapter.

For an $(\infty,2)$-category $\mathcal X$:
\begin{enumerate}[label=(\roman*)]
    \item $\mathcal X^{\op}$ denotes reversal of $1$-morphisms while
    retaining the directions of $2$-morphisms;

    \item $\mathcal X^{\mathrm{co}}$ denotes reversal of $2$-morphisms
    while retaining the directions of $1$-morphisms;

    \item $\mathcal X^{\mathrm{coop}}$ denotes simultaneous reversal of
    $1$- and $2$-morphisms.
\end{enumerate}
\end{convention}

\begin{convention}[Pseudonaturality orientation]
\label{conv:morita-pseudonaturality-orientation}
Let
\[
    F,G
    \colon
    \mathcal X
    \longrightarrow
    \mathcal Y
\]
be $(\infty,2)$-functors. A pseudonatural transformation
\[
    \eta
    \colon
    F
    \Longrightarrow
    G
\]
has, for every object
\[
    x\in\mathcal X,
\]
a $1$-morphism
\[
    \eta_x
    \colon
    F(x)
    \longrightarrow
    G(x).
\]

For a $1$-morphism
\[
    f
    \colon
    x
    \longrightarrow
    y,
\]
its pseudonaturality constraint is oriented as an invertible
$2$-morphism
\[
    G(f)\circ\eta_x
    \xRightarrow{\;\simeq\;}
    \eta_y\circ F(f).
\]

In a Morita category, horizontal composition is written from left to
right by relative tensor product. Thus the same pseudonaturality
constraint is written
\[
    \eta_x
    \otimes_{G(x)}
    G(f)
    \xRightarrow{\;\simeq\;}
    F(f)
    \otimes_{F(y)}
    \eta_y.
\]

All pseudonaturality constraints in this chapter use this orientation.
\end{convention}

\begin{definition}[Coherent weak duality involution]
\label{def:coherent-weak-duality-involution}
Let
\[
    \Cat_{(\infty,2)}^\simeq
\]
denote the maximal $\infty$-groupoid of $(\infty,2)$-categories and
equivalences in the fixed universe. The operation
\[
    \mathcal X
    \longmapsto
    \mathcal X^{\mathrm{co}}
\]
defines a coherent
\[
    C_2
\]
-action on
\[
    \Cat_{(\infty,2)}^\simeq.
\]

A \emph{coherent weak duality involution} on an $(\infty,2)$-category
\[
    \mathcal X
\]
is a coherent homotopy-fixed-point structure for this action whose
underlying object is
\[
    \mathcal X.
\]

Its first visible layers consist of:
\begin{enumerate}[label=(\roman*)]
    \item an equivalence of $(\infty,2)$-categories
    \[
        D
        \colon
        \mathcal X^{\mathrm{co}}
        \xrightarrow{\;\simeq\;}
        \mathcal X;
    \]

    \item a pseudonatural adjoint equivalence
    \[
        \operatorname{bid}^{D}
        \colon
        \id_{\mathcal X}
        \xRightarrow{\;\simeq\;}
        D D^{\mathrm{co}};
    \]

    \item an invertible modification
    \[
        \vartheta^{D}
        \colon
        D\ast
        \bigl(
            \operatorname{bid}^{D}
        \bigr)^{\mathrm{co}}
        \xRightarrow{\;\simeq\;}
        \operatorname{bid}^{D}\ast D
    \]
    between the two pseudonatural transformations
    \[
        D
        \Longrightarrow
        D D^{\mathrm{co}}D;
    \]

    \item the complete coherent
    \[
        C_2
    \]
    -cocycle data extending
    \[
        \operatorname{bid}^{D}
        \qquad\text{and}\qquad
        \vartheta^{D}.
    \]
\end{enumerate}

Here
\[
    \ast
\]
denotes pseudonatural whiskering. The homotopy-fixed-point formulation is
part of the definition. Thus all higher coherence beyond the displayed
layers is included rather than being inferred from pointwise
equivalences.
\end{definition}

\begin{definition}[Coherent $1$-op involution]
\label{def:coherent-one-op-involution}
Let
\[
    \Cat_{(\infty,2)}^\simeq
\]
be as in
\Ref{def:coherent-weak-duality-involution}. The operation
\[
    \mathcal X
    \longmapsto
    \mathcal X^{\op}
\]
defines a coherent
\[
    C_2
\]
-action on
\[
    \Cat_{(\infty,2)}^\simeq.
\]

A \emph{coherent $1$-op involution} on an $(\infty,2)$-category
\[
    \mathcal X
\]
is a coherent homotopy-fixed-point structure for this action whose
underlying object is
\[
    \mathcal X.
\]

Its first visible layers consist of:
\begin{enumerate}[label=(\roman*)]
    \item an equivalence
    \[
        O
        \colon
        \mathcal X^{\op}
        \xrightarrow{\;\simeq\;}
        \mathcal X;
    \]

    \item a pseudonatural adjoint equivalence
    \[
        \operatorname{bid}^{O}
        \colon
        \id_{\mathcal X}
        \xRightarrow{\;\simeq\;}
        O O^{\op};
    \]

    \item an invertible modification
    \[
        \vartheta^{O}
        \colon
        \left(
            O\ast
            \bigl(
                \operatorname{bid}^{O}
            \bigr)^{\op}
        \right)
        \circ
        \left(
            \operatorname{bid}^{O}\ast O
        \right)
        \xRightarrow{\;\simeq\;}
        \id_O.
    \]
    Explicitly,
    \[
        \operatorname{bid}^{O}\ast O
        \colon
        O
        \Longrightarrow
        O O^{\op}O,
    \]
    while
    \[
        O\ast
        \bigl(
            \operatorname{bid}^{O}
        \bigr)^{\op}
        \colon
        O O^{\op}O
        \Longrightarrow
        O;
    \]

    \item the complete coherent
    \[
        C_2
    \]
    -cocycle data extending
    \[
        \operatorname{bid}^{O}
        \qquad\text{and}\qquad
        \vartheta^{O}.
    \]
\end{enumerate}

The homotopy-fixed-point formulation is part of the definition. In
particular, the displayed triangle modification is only the first
visible involutivity coherence.
\end{definition}

\begin{convention}[Full dualizability]
\label{conv:morita-full-dualizability}
In a symmetric monoidal $(\infty,2)$-category, an object is
\emph{fully dualizable} if it is dualizable and its evaluation and
coevaluation $1$-morphisms admit both left and right adjoints.
\end{convention}

\section{Associative algebras and the higher Morita category}
\label{sec:associative-algebras-relative-base}

\begin{definition}[Associative algebra over \(A\)]
\label{def:associative-algebra-over-base}
An \emph{associative algebra over $A$} is an object
\[
    B
    \in
    \Alg_{\EE_1}(\mathcal W).
\]
We abbreviate
\[
    \Alg_A
    :=
    \Alg_{\EE_1}(\mathcal W).
\]
\end{definition}

\begin{remark}
\label{rem:commutative-base-noncommutative-algebra}
The algebra $A$ is commutative because it defines the relative affine
base. The algebra $B$ is only associative and need not be commutative.
\end{remark}

\begin{definition}[Opposite and enveloping algebra]
\label{def:opposite-enveloping-algebra}
For
\[
    B\in\Alg_A,
\]
let
\[
    B^{\op}
\]
denote its opposite algebra and define
\[
    B^e
    :=
    B\otimes_A B^{\op}.
\]
\end{definition}

\begin{definition}[Bimodule]
\label{def:relative-bimodule}
For
\[
    B,C\in\Alg_A,
\]
define
\[
    \BMod_{B,C}(\mathcal W)
    :=
    \LMod_{B\otimes_A C^{\op}}(\mathcal W).
\]
An object
\[
    {}_B M_C
    \in
    \BMod_{B,C}(\mathcal W)
\]
is called a \emph{$B$--$C$-bimodule}. It is regarded as a Morita morphism
\[
    M
    \colon
    B
    \longrightarrow
    C.
\]
\end{definition}

\begin{construction}[Morita composition]
\label{con:morita-composition}
For bimodules
\[
    M
    \colon
    B
    \longrightarrow
    C,
    \qquad
    N
    \colon
    C
    \longrightarrow
    D,
\]
their Morita composite is
\[
    M\otimes_C N,
\]
regarded as a $B$--$D$-bimodule.
\end{construction}

\begin{theorem}[Higher Morita category]
\label{thm:higher-morita-category}
There is an $(\infty,2)$-category
\[
    \Mor_A
    :=
    \Alg_{\EE_1}^{\mathrm{Mor}}(\mathcal W)
\]
whose:
\begin{enumerate}[label=(\roman*)]
    \item objects are associative $A$-algebras;

    \item $1$-morphisms are bimodules;

    \item $2$-morphisms are bimodule maps;

    \item horizontal composition is relative tensor product.
\end{enumerate}
\end{theorem}

\begin{proof}
Because
\[
    \mathcal W
\]
is presentable and its tensor product preserves colimits separately in
each variable, all two-sided bar constructions defining relative tensor
products exist. Their geometric realizations are preserved by tensoring
on either side.

Thus
\[
    \mathcal W
\]
has good relative tensor products in the sense required by the higher
Morita construction. The result follows from
\cite[Theorem~1.1]{HaugsengHigherMorita}.
\end{proof}

\begin{proposition}[Symmetric monoidal higher Morita category]
\label{prop:symmetric-monoidal-higher-morita-category}
The symmetric monoidal structure on
\[
    \mathcal W
    =
    \Mod_A(\V)
\]
induces a canonical symmetric monoidal structure
\[
    \Mor_A^\otimes
\]
on the higher Morita $(\infty,2)$-category.

Its tensor unit is the commutative base algebra
\[
    A.
\]
On objects, the tensor product is
\[
    (B,C)
    \longmapsto
    B\otimes_A C.
\]

For bimodules
\[
    M
    \colon
    B
    \longrightarrow
    C,
    \qquad
    M'
    \colon
    B'
    \longrightarrow
    C',
\]
their external tensor product is
\[
    M\otimes_A M'
    \colon
    B\otimes_A B'
    \longrightarrow
    C\otimes_A C'.
\]
The tensor product of $2$-morphisms is the tensor product of their
underlying bimodule maps.

For composable bimodules
\[
    M
    \colon
    B
    \longrightarrow
    C,
    \qquad
    N
    \colon
    C
    \longrightarrow
    D,
\]
and
\[
    M'
    \colon
    B'
    \longrightarrow
    C',
    \qquad
    N'
    \colon
    C'
    \longrightarrow
    D',
\]
there is a canonical interchange equivalence
\[
\begin{aligned}
    &
    (M\otimes_A M')
    \otimes_{C\otimes_A C'}
    (N\otimes_A N')
    \\
    &\qquad\xrightarrow{\;\simeq\;}
    (M\otimes_C N)
    \otimes_A
    (M'\otimes_{C'}N').
\end{aligned}
\]
These equivalences satisfy all symmetric monoidal and interchange
coherences.
\end{proposition}

\begin{proof}
The external tensor product carries a $B$--$C$-bimodule and a
$B'$--$C'$-bimodule to a
\[
    (B\otimes_A B')
    \text{--}
    (C\otimes_A C')
\]
-bimodule. It is functorial on bimodule maps.

The interchange equivalence is induced by the canonical comparison
between the corresponding iterated two-sided bar constructions:
\[
\begin{aligned}
    &
    \left|
        \operatorname{Bar}
        \bigl(
            M\otimes_A M',
            C\otimes_A C',
            N\otimes_A N'
        \bigr)
    \right|
    \\
    &\qquad\xrightarrow{\;\simeq\;}
    \left|
        \operatorname{Bar}(M,C,N)
    \right|
    \otimes_A
    \left|
        \operatorname{Bar}(M',C',N')
    \right|.
\end{aligned}
\]
Because tensor product in
\[
    \mathcal W
\]
preserves geometric realizations separately in each variable, the
comparison is an equivalence. Its compatibility with iterated
composition follows from the corresponding comparison of
multisimplicial bar constructions.

The unit constraints are induced by
\[
    A\otimes_A B
    \simeq
    B
    \simeq
    B\otimes_A A
\]
and their analogues for bimodules.

More conceptually,
\cite[Corollary~5.42]{HaugsengHigherMorita}
constructs the higher Morita assignment as a lax symmetric monoidal
functor on monoidal $\infty$-categories with good relative tensor
products and monoidal functors preserving them.

The symmetric monoidal $\infty$-category
\[
    \mathcal W
\]
has good relative tensor products by the hypotheses used in
\Ref{thm:higher-morita-category}. Applying Haugseng's lax symmetric
monoidal Morita construction to the commutative algebra object
\[
    \mathcal W
\]
in monoidal $\infty$-categories therefore produces a commutative algebra
object
\[
    \Mor_A^\otimes
\]
in $(\infty,2)$-categories.

Unwinding this construction gives the tensor products, unit, and
interchange equivalences displayed above. Since they arise from a single
lax symmetric monoidal construction, all associativity, unit, symmetry,
and interchange coherences are supplied simultaneously.
\end{proof}

\section{Internal Homs and one-sided perfectness}
\label{sec:right-perfect-morita-adjoints}

For an associative algebra
\[
    C,
\]
the category of right $C$-modules is not generally monoidal.
Consequently, right-perfectness cannot be defined as ordinary monoidal
dualizability inside
\[
    \RMod_C(\mathcal W).
\]
It is defined using the Morita morphism from the commutative base algebra
\[
    A
\]
to
\[
    C.
\]

\subsection{Internal Homs of right modules}

\begin{proposition}[Existence of right-module internal Homs]
\label{prop:internal-hom-right-modules-exists}
Let
\[
    M\in\RMod_C(\mathcal W).
\]
The functor
\[
    -\otimes_A M
    \colon
    \mathcal W
    \longrightarrow
    \RMod_C(\mathcal W)
\]
admits a right adjoint.
\end{proposition}

\begin{proof}
The displayed functor is accessible and preserves all colimits. Both its
source and target are presentable $\infty$-categories. The adjoint
functor theorem therefore supplies a right adjoint.
\end{proof}

\begin{definition}[Internal Hom of right modules]
\label{def:internal-hom-right-modules}
For
\[
    M,N\in\RMod_C(\mathcal W),
\]
the object
\[
    \underline{\Hom}_C(M,N)
    \in
    \mathcal W
\]
is defined by the natural equivalences
\[
    \Map_{\RMod_C(\mathcal W)}
    \bigl(
        X\otimes_A M,
        N
    \bigr)
    \simeq
    \Map_{\mathcal W}
    \bigl(
        X,
        \underline{\Hom}_C(M,N)
    \bigr)
\]
for every
\[
    X\in\mathcal W.
\]
\end{definition}

\begin{definition}[Right \(C\)-dualizable module]
\label{def:right-c-dualizable-module}
A right $C$-module
\[
    M
\]
is \emph{right $C$-dualizable} if, regarded as an
$A$--$C$-bimodule
\[
    M
    \colon
    A
    \longrightarrow
    C,
\]
it admits a right adjoint in the Morita $(\infty,2)$-category.

Explicitly, this consists of a $C$--$A$-bimodule
\[
    M_C^\vee
    \colon
    C
    \longrightarrow
    A
\]
and morphisms
\[
    \eta_M^{A}
    \colon
    A
    \longrightarrow
    M\otimes_C M_C^\vee,
\]
\[
    \epsilon_M^{A}
    \colon
    M_C^\vee\otimes_A M
    \longrightarrow
    C
\]
satisfying the two triangle identities.
\end{definition}

\begin{proposition}[Equivariant tensor--Hom formula]
\label{prop:right-module-tensor-hom}
Let
\[
    M
    \in
    \BMod_{B,C}(\mathcal W)
\]
be right $C$-dualizable as an $A$--$C$-bimodule, and let
\[
    N
    \in
    \BMod_{D,C}(\mathcal W).
\]
Then:
\begin{enumerate}[label=(\roman*)]
    \item the right dual
    \[
        M_C^\vee
    \]
    has a natural $C$--$B$-bimodule structure;

    \item the object
    \[
        \underline{\Hom}_C(M,N)
    \]
    has a natural $D$--$B$-bimodule structure;

    \item there is a canonical equivalence
    \[
        N\otimes_C M_C^\vee
        \xrightarrow{\;\simeq\;}
        \underline{\Hom}_C(M,N)
    \]
    in
    \[
        \BMod_{D,B}(\mathcal W).
    \]
\end{enumerate}

In particular,
\[
    M_C^\vee
    \simeq
    \underline{\Hom}_C(M,C)
\]
as a $C$--$B$-bimodule.
\end{proposition}

\begin{proof}
The left $B$-action on
\[
    M
\]
induces a right $B$-action on
\[
    M_C^\vee
\]
by precomposition. Thus
\[
    M_C^\vee
\]
is a $C$--$B$-bimodule.

The left $D$-action on
\[
    N
\]
induces a left $D$-action on
\[
    \underline{\Hom}_C(M,N),
\]
while precomposition with the left $B$-action on
\[
    M
\]
induces a right $B$-action. Thus
\[
    \underline{\Hom}_C(M,N)
\]
is a $D$--$B$-bimodule.

Define
\[
    N\otimes_C M_C^\vee
    \longrightarrow
    \underline{\Hom}_C(M,N)
\]
as the adjoint of
\[
\begin{aligned}
    (N\otimes_C M_C^\vee)\otimes_A M
    &\simeq
    N\otimes_C
    (M_C^\vee\otimes_A M)
    \\
    &\xrightarrow{
        \id_N\otimes\epsilon_M^{A}
    }
    N\otimes_C C
    \\
    &\simeq
    N.
\end{aligned}
\]

For every
\[
    X\in\mathcal W,
\]
the adjunction
\[
    M\dashv M_C^\vee
\]
gives
\[
\begin{aligned}
    \Map_{\mathcal W}
    \bigl(
        X,
        N\otimes_C M_C^\vee
    \bigr)
    &\simeq
    \Map_{\RMod_C(\mathcal W)}
    \bigl(
        X\otimes_A M,
        N
    \bigr)
    \\
    &\simeq
    \Map_{\mathcal W}
    \bigl(
        X,
        \underline{\Hom}_C(M,N)
    \bigr).
\end{aligned}
\]
Yoneda identifies the two underlying objects.

Naturality with respect to endomorphisms of
\[
    N
\]
gives compatibility with the left $D$-action. Naturality
contravariantly with respect to endomorphisms of
\[
    M
\]
gives compatibility with the right $B$-action. Hence the equivalence is
an equivalence of $D$--$B$-bimodules.

Taking
\[
    N=C
\]
gives
\[
    M_C^\vee
    \simeq
    \underline{\Hom}_C(M,C).
\]
\end{proof}

\subsection{Right-perfect and left-perfect bimodules}

\begin{definition}[Right-perfect bimodule]
\label{def:right-perfect-bimodule}
A $B$--$C$-bimodule
\[
    M
    \colon
    B
    \longrightarrow
    C
\]
is \emph{right perfect} if its underlying $A$--$C$-bimodule is right
$C$-dualizable.
\end{definition}

\begin{definition}[Action reversal]
\label{def:action-reversal-bimodule}
For a $B$--$C$-bimodule
\[
    M
    \colon
    B
    \longrightarrow
    C,
\]
let
\[
    M^{\op}
    \colon
    C^{\op}
    \longrightarrow
    B^{\op}
\]
denote the same underlying object with its left and right actions
reversed.
\end{definition}

\begin{definition}[Left-perfect bimodule]
\label{def:left-perfect-bimodule}
A $B$--$C$-bimodule
\[
    M
\]
is \emph{left perfect} if the action-reversed bimodule
\[
    M^{\op}
    \colon
    C^{\op}
    \longrightarrow
    B^{\op}
\]
is right perfect.
\end{definition}

\section{Right-perfectness and Morita adjointability}
\label{sec:right-perfectness-morita-adjointability}

For every associative $A$-algebra
\[
    B,
\]
let
\[
    \iota_B
    :=
    {}_A B_B
    \colon
    A
    \longrightarrow
    B
\]
denote the regular extension bimodule.

\begin{lemma}[The regular extension adjunction]
\label{lem:regular-extension-adjunction}
The bimodule
\[
    \iota_B
    =
    {}_A B_B
\]
has right adjoint
\[
    \iota_B^R
    :=
    {}_B B_A.
\]
\end{lemma}

\begin{proof}
The unit is the algebra unit
\[
    A
    \longrightarrow
    B
    \simeq
    B\otimes_B B,
\]
regarded as a morphism
\[
    A
    \longrightarrow
    \iota_B\otimes_B\iota_B^R.
\]
The counit is multiplication
\[
    B\otimes_A B
    \longrightarrow
    B,
\]
regarded as a morphism
\[
    \iota_B^R\otimes_A\iota_B
    \longrightarrow
    B.
\]

The first triangle composite is
\[
\begin{aligned}
    B
    &\simeq
    A\otimes_A B
    \\
    &\longrightarrow
    B\otimes_A B
    \\
    &\longrightarrow
    B,
\end{aligned}
\]
which is the left unit law.

The second triangle composite is
\[
\begin{aligned}
    B
    &\simeq
    B\otimes_A A
    \\
    &\longrightarrow
    B\otimes_A B
    \\
    &\longrightarrow
    B,
\end{aligned}
\]
which is the right unit law.
\end{proof}

\begin{proposition}[Right-perfectness equals Morita right-adjointability]
\label{prop:right-perfect-morita-adjointability}
Let
\[
    M
    \colon
    B
    \longrightarrow
    C
\]
be a $B$--$C$-bimodule. The following conditions are equivalent:
\begin{enumerate}[label=(\roman*)]
    \item $M$ is right perfect;

    \item $M$ admits a right adjoint as a $1$-morphism
    \[
        B
        \longrightarrow
        C
    \]
    in
    \[
        \Mor_A.
    \]
\end{enumerate}

When these conditions hold, the right adjoint is canonically equivalent
to
\[
    M^R
    :=
    \underline{\Hom}_C(M,C)
    \colon
    C
    \longrightarrow
    B.
\]
\end{proposition}
\begin{proof}
Suppose first that
\[
    M
\]
is right perfect. Its underlying $A$--$C$-bimodule admits a right
adjoint
\[
    M_C^\vee
    \colon
    C
    \longrightarrow
    A
\]
with unit and counit
\[
    \eta_M^{A}
    \colon
    A
    \longrightarrow
    M\otimes_C M_C^\vee,
\]
\[
    \epsilon_M^{A}
    \colon
    M_C^\vee\otimes_A M
    \longrightarrow
    C.
\]

By
\Ref{prop:right-module-tensor-hom},
the object
\[
    M^R
    :=
    \underline{\Hom}_C(M,C)
\]
has a canonical $C$--$B$-bimodule structure, and there is a canonical
equivalence
\[
    M^R
    \simeq
    M_C^\vee
\]
of $C$--$B$-bimodules. We use this equivalence to identify the two
objects.

Applying
\Ref{prop:right-module-tensor-hom}
with
\[
    N=M
\]
gives a canonical equivalence of $B$--$B$-bimodules
\[
    \theta_M
    \colon
    M\otimes_C M^R
    \xrightarrow{\;\simeq\;}
    \underline{\Hom}_C(M,M).
\]

Under the tensor--Hom adjunction, the left $B$-action
\[
    B\otimes_A M
    \longrightarrow
    M
\]
has an adjoint transpose
\[
    B
    \longrightarrow
    \underline{\Hom}_C(M,M).
\]
Transporting this morphism through
\[
    \theta_M^{-1}
\]
defines
\[
    \eta_M^{B}
    \colon
    B
    \longrightarrow
    M\otimes_C M^R.
\]

The composite
\[
    A
    \longrightarrow
    B
    \xrightarrow{\eta_M^{B}}
    M\otimes_C M^R
\]
is canonically equivalent to
\[
    \eta_M^{A}.
\]
Indeed, under
\[
    \theta_M,
\]
both morphisms correspond to the identity endomorphism
\[
    \id_M.
\]

Evaluation gives an $A$-balanced morphism
\[
    M^R\otimes_A M
    \longrightarrow
    C.
\]
It is $B$-balanced because
\[
    \ev\bigl((\phi\cdot b)\otimes m\bigr)
    \simeq
    \phi(bm)
    \simeq
    \ev(\phi\otimes bm).
\]
It therefore descends uniquely to a morphism
\[
    \epsilon_M^{C}
    \colon
    M^R\otimes_B M
    \longrightarrow
    C.
\]

Let
\[
    q_M
    \colon
    M^R\otimes_A M
    \longrightarrow
    M^R\otimes_B M
\]
denote the canonical balancing morphism. By construction,
\[
    \epsilon_M^{A}
    \simeq
    \epsilon_M^{C}\circ q_M.
\]

We now verify the triangle identities.

The first triangle composite is
\[
\begin{aligned}
    M
    &\simeq
    B\otimes_B M
    \\
    &\xrightarrow{
        \eta_M^{B}\otimes\id_M
    }
    M\otimes_C M^R\otimes_B M
    \\
    &\xrightarrow{
        \id_M\otimes\epsilon_M^{C}
    }
    M.
\end{aligned}
\]
Under the equivalence
\[
    M\otimes_C M^R
    \simeq
    \underline{\Hom}_C(M,M),
\]
the morphism
\[
    \eta_M^{B}
\]
is the adjoint transpose of the left $B$-action. Hence the displayed
composite is evaluation of that action on the regular element of
\[
    B.
\]
Equivalently, it is evaluation of
\[
    \id_M.
\]
It is therefore canonically equivalent to
\[
    \id_M.
\]

For the second triangle, consider
\[
\begin{aligned}
    M^R
    &\simeq
    M^R\otimes_B B
    \\
    &\xrightarrow{
        \id_{M^R}\otimes\eta_M^{B}
    }
    M^R\otimes_B M\otimes_C M^R
    \\
    &\xrightarrow{
        \epsilon_M^{C}\otimes\id_{M^R}
    }
    M^R.
\end{aligned}
\]
There is a canonically commutative diagram
\[
\begin{tikzcd}[column sep=large]
M^R\otimes_A A
    \ar[r, "\id_{M^R}\otimes\eta_M^{A}"]
    \ar[d, "\simeq"']
&
M^R\otimes_A M\otimes_C M^R
    \ar[d, "q_M\otimes\id_{M^R}"]
\\
M^R\otimes_B B
    \ar[r, "\id_{M^R}\otimes\eta_M^{B}"']
&
M^R\otimes_B M\otimes_C M^R.
\end{tikzcd}
\]
The left vertical equivalence is induced by the right $B$-action on
\[
    M^R,
\]
and commutativity follows from
\[
    \eta_M^{B}|_A
    \simeq
    \eta_M^{A}.
\]

Moreover,
\[
    \epsilon_M^{A}
    \simeq
    \epsilon_M^{C}\circ q_M.
\]
Consequently, the second Morita triangle composite is canonically
identified with
\[
\begin{aligned}
    M^R
    &\simeq
    M^R\otimes_A A
    \\
    &\xrightarrow{
        \id_{M^R}\otimes\eta_M^{A}
    }
    M^R\otimes_A M\otimes_C M^R
    \\
    &\xrightarrow{
        \epsilon_M^{A}\otimes\id_{M^R}
    }
    M^R.
\end{aligned}
\]
This is the second triangle composite for the underlying adjunction
\[
    M\dashv M_C^\vee.
\]
It is therefore canonically equivalent to
\[
    \id_{M^R}.
\]

Thus
\[
    M\dashv M^R
\]
in
\[
    \Mor_A.
\]

Conversely, suppose that
\[
    M
    \colon
    B
    \longrightarrow
    C
\]
admits a right adjoint in
\[
    \Mor_A.
\]
Its underlying $A$--$C$-bimodule is the Morita composite
\[
    A
    \xrightarrow{\iota_B}
    B
    \xrightarrow{M}
    C,
\]
because
\[
    \iota_B\otimes_B M
    \simeq
    B\otimes_B M
    \simeq
    M.
\]

By
\Ref{lem:regular-extension-adjunction},
the morphism
\[
    \iota_B
\]
admits a right adjoint. By assumption,
\[
    M
\]
admits a right adjoint. Pasting these adjunctions therefore gives a
right adjoint for their composite
\[
    \iota_B\otimes_B M.
\]
Hence the underlying $A$--$C$-bimodule of
\[
    M
\]
is right $C$-dualizable. Thus
\[
    M
\]
is right perfect.

Finally, the right adjoint of
\[
    M
\]
is canonically equivalent to
\[
    \underline{\Hom}_C(M,C)
\]
by
\Ref{prop:right-module-tensor-hom}
and the contractible uniqueness of right adjoints.
\end{proof}

\begin{corollary}[Right Morita adjoint]
\label{cor:right-morita-adjoint}
Let
\[
    M
    \colon
    B
    \longrightarrow
    C
\]
be right perfect. Then
\[
    M
\]
admits a right adjoint
\[
    M^R
    :=
    \underline{\Hom}_C(M,C)
    \colon
    C
    \longrightarrow
    B.
\]
\end{corollary}

\begin{proof}
This is the forward implication of
\Ref{prop:right-perfect-morita-adjointability}.
\end{proof}

\begin{proposition}[Reversal of Morita composition]
\label{prop:op-reversal-morita-composition}
For composable bimodules
\[
    M
    \colon
    B
    \longrightarrow
    C,
    \qquad
    N
    \colon
    C
    \longrightarrow
    D,
\]
there is a canonical equivalence
\[
    (M\otimes_C N)^{\op}
    \xrightarrow{\;\simeq\;}
    N^{\op}
    \otimes_{C^{\op}}
    M^{\op}.
\]
These equivalences are coherently associative and unital.
\end{proposition}

\begin{proof}
The relative tensor product
\[
    M\otimes_C N
\]
is the geometric realization of the two-sided bar construction
\[
    [n]
    \longmapsto
    M
    \otimes_A
    C^{\otimes_A n}
    \otimes_A
    N.
\]

After action reversal and reversal of the order of tensor factors, this
simplicial object identifies with
\[
    [n]
    \longmapsto
    N^{\op}
    \otimes_A
    (C^{\op})^{\otimes_A n}
    \otimes_A
    M^{\op},
\]
which is the two-sided bar construction defining
\[
    N^{\op}
    \otimes_{C^{\op}}
    M^{\op}.
\]

The comparison is natural and is induced levelwise by the symmetry of
\[
    \mathcal W.
\]
Its associativity and unit coherences follow because all comparisons are
induced by the same coherent reversal of the corresponding iterated bar
constructions.
\end{proof}

\begin{corollary}[Left-perfectness equals Morita left-adjointability]
\label{cor:left-perfect-morita-adjointability}
Let
\[
    M
    \colon
    B
    \longrightarrow
    C
\]
be a bimodule. The following conditions are equivalent:
\begin{enumerate}[label=(\roman*)]
    \item $M$ is left perfect;

    \item $M$ admits a left adjoint in
    \[
        \Mor_A.
    \]
\end{enumerate}
\end{corollary}

\begin{proof}
By definition,
\[
    M
\]
is left perfect if and only if
\[
    M^{\op}
\]
is right perfect. By
\Ref{prop:right-perfect-morita-adjointability},
this is equivalent to
\[
    M^{\op}
\]
admitting a right adjoint.

Suppose
\[
    M^{\op}\dashv N.
\]
Action reversal transforms this adjunction into
\[
    N^{\op}\dashv M.
\]
Thus
\[
    M
\]
admits a left adjoint.

Conversely, action reversal of a left adjunction for
\[
    M
\]
produces a right adjunction for
\[
    M^{\op}.
\]
\end{proof}

\begin{proposition}[Adjoint of a composite]
\label{prop:morita-adjoint-composite}
Let
\[
    M
    \colon
    B
    \longrightarrow
    C,
    \qquad
    N
    \colon
    C
    \longrightarrow
    D
\]
be right perfect. Then
\[
    M\otimes_C N
\]
is right perfect and
\[
    (M\otimes_C N)^R
    \simeq
    N^R\otimes_C M^R.
\]
\end{proposition}

\begin{proof}
By
\Ref{prop:right-perfect-morita-adjointability},
the bimodules
\[
    M
    \qquad\text{and}\qquad
    N
\]
are right-adjointable. Pasting their adjunctions gives an adjunction
\[
    M\otimes_C N
    \dashv
    N^R\otimes_C M^R.
\]

The unit is the composite
\[
\begin{aligned}
    B
    &\xrightarrow{\eta_M^{B}}
    M\otimes_C M^R
    \\
    &\xrightarrow{
        \id_M
        \otimes
        \eta_N^{C}
        \otimes
        \id_{M^R}
    }
    M\otimes_C N\otimes_D N^R\otimes_C M^R.
\end{aligned}
\]
The counit is
\[
\begin{aligned}
    N^R\otimes_C M^R
    \otimes_B
    M\otimes_C N
    &\xrightarrow{
        \id_{N^R}
        \otimes
        \epsilon_M^{C}
        \otimes
        \id_N
    }
    N^R\otimes_C N
    \\
    &\xrightarrow{\epsilon_N^{D}}
    D.
\end{aligned}
\]

The triangle identities follow by pasting the triangle identities for
\[
    M\dashv M^R
    \qquad\text{and}\qquad
    N\dashv N^R.
\]
Thus the composite is right-adjointable and hence right perfect.
The coherent compatibility of these adjoints with iterated composition is
proved in
\Ref{prop:coherent-adjoints-composites}.
\end{proof}

\section{Smoothness, properness, and full dualizability}
\label{sec:smooth-proper-full-dualizable}

\begin{definition}[Proper and smooth algebra]
\label{def:smooth-proper-algebra}
Let
\[
    B
    \in
    \Alg_A
\]
and let
\[
    B^e
    :=
    B\otimes_A B^{\op}.
\]

The algebra
\[
    B
\]
is:
\begin{enumerate}[label=(\roman*)]
    \item \emph{proper over $A$} if its underlying $A$-module belongs to
    \[
        \Perf(A);
    \]

    \item \emph{smooth over $A$} if the evaluation, or diagonal,
    bimodule
    \[
        {}_{B^e}B_A
        \colon
        B^e
        \longrightarrow
        A
    \]
    is left perfect.
\end{enumerate}

Equivalently,
\[
    B
\]
is smooth if the action-reversed diagonal bimodule
\[
    {}_A B_{(B^e)^{\op}}
    \colon
    A
    \longrightarrow
    (B^e)^{\op}
\]
is right perfect.

Equivalently again, the diagonal object
\[
    {}_{B^e}B
\]
is one-sided perfect as a left $B^e$-module in the Morita sense: under
action reversal it determines a right-perfect Morita morphism
\[
    A
    \longrightarrow
    (B^e)^{\op}.
\]
\end{definition}

\begin{remark}[Comparison with compact-module terminology]
\label{rem:smoothness-compact-module-terminology}
In settings in which the right-perfect modules over an associative
algebra agree with the compact modules, the smoothness condition says
that
\[
    {}_{B^e}B
\]
is compact, or equivalently perfect in the usual thick-subcategory
sense, as a left $B^e$-module. No such compact-generation identification
is used in this chapter.
\end{remark}

\begin{theorem}[Smooth, proper, and fully dualizable]
\label{thm:smooth-proper-fully-dualizable}
Every associative $A$-algebra
\[
    B
\]
is dualizable as an object of the symmetric monoidal Morita
$(\infty,2)$-category
\[
    \Mor_A^\otimes.
\]
Its dual is
\[
    B^{\op}.
\]

Moreover,
\[
    B
\]
is fully dualizable in
\[
    \Mor_A^\otimes
\]
if and only if it is smooth and proper over
\[
    A.
\]
\end{theorem}

\begin{proof}
The tensor dual of
\[
    B
\]
is
\[
    B^{\op}.
\]
The evaluation and coevaluation $1$-morphisms are the diagonal bimodule
with the appropriate action conventions:
\[
    \ev_B
    :=
    {}_{B\otimes_A B^{\op}}B_A
    \colon
    B\otimes_A B^{\op}
    \longrightarrow
    A
\]
and
\[
    \coev_B
    :=
    {}_A B_{B^{\op}\otimes_A B}
    \colon
    A
    \longrightarrow
    B^{\op}\otimes_A B.
\]

The first duality triangle is the composite $B$--$B$-bimodule
\[
\begin{aligned}
    &\bigl(
        \id_B\otimes_A\coev_B
    \bigr)
    \otimes_{B\otimes_A B^{\op}\otimes_A B}
    \bigl(
        \ev_B\otimes_A\id_B
    \bigr).
\end{aligned}
\]
Multiplication in
\[
    B
\]
induces a canonical equivalence
\[
\begin{aligned}
    &\bigl(
        \id_B\otimes_A\coev_B
    \bigr)
    \otimes_{B\otimes_A B^{\op}\otimes_A B}
    \bigl(
        \ev_B\otimes_A\id_B
    \bigr)
    \xrightarrow{\;\simeq\;}
    {}_B B_B.
\end{aligned}
\]
In the two-sided bar model, this equivalence is the regular Morita
contraction; the algebra unit supplies a contracting extra degeneracy.

Similarly, the second duality triangle is the composite
$B^{\op}$--$B^{\op}$-bimodule
\[
\begin{aligned}
    &\bigl(
        \coev_B\otimes_A\id_{B^{\op}}
    \bigr)
    \otimes_{B^{\op}\otimes_A B\otimes_A B^{\op}}
    \bigl(
        \id_{B^{\op}}\otimes_A\ev_B
    \bigr),
\end{aligned}
\]
and multiplication induces a canonical equivalence
\[
\begin{aligned}
    &\bigl(
        \coev_B\otimes_A\id_{B^{\op}}
    \bigr)
    \otimes_{B^{\op}\otimes_A B\otimes_A B^{\op}}
    \bigl(
        \id_{B^{\op}}\otimes_A\ev_B
    \bigr)
    \xrightarrow{\;\simeq\;}
    {}_{B^{\op}}B^{\op}_{B^{\op}}.
\end{aligned}
\]
The two equivalences satisfy the triangle coherences because both arise
from the unit and associativity coherences of the regular bimodule.
Thus
\[
    B
\]
is dualizable with dual
\[
    B^{\op}.
\]

It remains to determine when the evaluation and coevaluation
$1$-morphisms admit both adjoints.

The evaluation bimodule is
\[
    \ev_B
    =
    {}_{B^e}B_A
    \colon
    B^e
    \longrightarrow
    A.
\]
By
\Ref{prop:right-perfect-morita-adjointability},
it has a right adjoint if and only if it is right perfect.

Since its target algebra is
\[
    A,
\]
this means precisely that its underlying $A$--$A$-bimodule
\[
    {}_A B_A
\]
is dualizable under
\[
    \otimes_A.
\]
The endomorphism monoidal category of
\[
    A
\]
in
\[
    \Mor_A
\]
is canonically
\[
    \Mod_A(\V).
\]
Therefore
\[
    \ev_B
    \text{ has a right adjoint}
    \quad\Longleftrightarrow\quad
    B\in\Perf(A).
\]
This is properness.

By
\Ref{cor:left-perfect-morita-adjointability},
the evaluation bimodule has a left adjoint if and only if it is left
perfect. This is precisely smoothness.

Action reversal identifies the evaluation with the coevaluation, up to
the canonical equivalence
\[
    (B\otimes_A B^{\op})^{\op}
    \simeq
    B^{\op}\otimes_A B.
\]
Action reversal interchanges left and right adjoints. Hence evaluation
and coevaluation both admit left and right adjoints exactly when
\[
    B
\]
is both smooth and proper.

Therefore
\[
    B
\]
is fully dualizable if and only if it is smooth and proper over
\[
    A.
\]
\end{proof}

\begin{remark}
\label{rem:smooth-proper-full-dualizable}
The smooth and proper algebras form the fully dualizable sector of the
Morita category. Compare
\cite{ToenDerivedMorita,HaugsengHigherMorita}.
\end{remark}

\section{The elementary mate correspondence}
\label{sec:elementary-mate-correspondence}

\begin{definition}[Right mate]
\label{def:right-mate}
Let
\[
    \alpha
    \colon
    M
    \longrightarrow
    N
\]
be a map between right-perfect $B$--$C$-bimodules. Its \emph{right mate}
is the map
\[
    \alpha^R
    \colon
    N^R
    \longrightarrow
    M^R
\]
defined by the composite
\[
\begin{aligned}
    N^R
    &\simeq
    N^R\otimes_B B
    \\
    &\xrightarrow{
        \id_{N^R}\otimes\eta_M^{B}
    }
    N^R\otimes_B M\otimes_C M^R
    \\
    &\xrightarrow{
        \id_{N^R}\otimes\alpha\otimes\id_{M^R}
    }
    N^R\otimes_B N\otimes_C M^R
    \\
    &\xrightarrow{
        \epsilon_N^{C}\otimes\id_{M^R}
    }
    M^R.
\end{aligned}
\]
\end{definition}

\begin{proposition}[Mate equivalence]
\label{prop:mate-equivalence}
For right-perfect $B$--$C$-bimodules
\[
    M
    \qquad\text{and}\qquad
    N,
\]
taking right mates defines an equivalence of mapping spaces
\[
    \Map_{\BMod_{B,C}(\mathcal W)}(M,N)
    \xrightarrow{\;\simeq\;}
    \Map_{\BMod_{C,B}(\mathcal W)}(N^R,M^R).
\]

Its inverse sends
\[
    \beta
    \colon
    N^R
    \longrightarrow
    M^R
\]
to the composite
\[
\begin{aligned}
    \beta^L
    \colon
    M
    &\simeq
    B\otimes_B M
    \\
    &\xrightarrow{
        \eta_N^{B}\otimes\id_M
    }
    N\otimes_C N^R\otimes_B M
    \\
    &\xrightarrow{
        \id_N\otimes\beta\otimes\id_M
    }
    N\otimes_C M^R\otimes_B M
    \\
    &\xrightarrow{
        \id_N\otimes\epsilon_M^{C}
    }
    N.
\end{aligned}
\]

Consequently,
\[
    (\id_M)^R
    \simeq
    \id_{M^R}
\]
and
\[
    (\beta\circ\alpha)^R
    \simeq
    \alpha^R\circ\beta^R.
\]
\end{proposition}

\begin{proof}
The right mate
\[
    \alpha^R
\]
is characterized by the counit equation
\[
    \epsilon_M^{C}
    \circ
    (\alpha^R\otimes_B\id_M)
    \simeq
    \epsilon_N^{C}
    \circ
    (\id_{N^R}\otimes_B\alpha)
\]
as maps
\[
    N^R\otimes_B M
    \longrightarrow
    C.
\]

Substituting the displayed formulas for
\[
    \alpha^R
    \qquad\text{and}\qquad
    \beta^L,
\]
the triangle identities for
\[
    M\dashv M^R
    \qquad\text{and}\qquad
    N\dashv N^R
\]
give canonical equivalences
\[
    (\alpha^R)^L
    \simeq
    \alpha,
    \qquad
    (\beta^L)^R
    \simeq
    \beta.
\]
Thus mating is an equivalence of mapping spaces.

The formulas for identities and composites follow either by direct
pasting or from the uniqueness characterization by the counit equation.
\end{proof}

\section{The coherent mate calculus}
\label{sec:coherent-mate-calculus}

Let
\[
    \mathcal X
\]
be an $(\infty,2)$-category. Write
\[
    \mathcal X^{\mathrm L}
\]
for the wide, locally full sub-$(\infty,2)$-category whose
$1$-morphisms admit right adjoints, and write
\[
    \mathcal X^{\mathrm R}
\]
for the wide, locally full sub-$(\infty,2)$-category whose
$1$-morphisms admit left adjoints.

Thus the superscripts record the side occupied by the selected
$1$-morphism: a morphism of
\[
    \mathcal X^{\mathrm L}
\]
is a left adjoint, and a morphism of
\[
    \mathcal X^{\mathrm R}
\]
is a right adjoint.

\begin{theorem}[Coherent mate calculus]
\label{thm:coherent-mate-calculus}
For every $(\infty,2)$-category
\[
    \mathcal X,
\]
taking right adjoints and right mates defines an equivalence of
$(\infty,2)$-categories
\[
    R_{\mathcal X}
    \colon
    \bigl(
        \mathcal X^{\mathrm L}
    \bigr)^{\mathrm{coop}}
    \xrightarrow{\;\simeq\;}
    \mathcal X^{\mathrm R}.
\]

On the underlying $1$-morphisms of
\[
    \mathcal X,
\]
it sends a right-adjointable $1$-morphism
\[
    f
    \colon
    x
    \longrightarrow
    y
\]
to a right adjoint
\[
    f^R
    \colon
    y
    \longrightarrow
    x,
\]
and a $2$-morphism
\[
    \alpha
    \colon
    f
    \Longrightarrow
    g
\]
to its right mate
\[
    \alpha^R
    \colon
    g^R
    \Longrightarrow
    f^R.
\]

Consequently:
\begin{enumerate}[label=(\roman*)]
    \item for composable right-adjointable morphisms
    \[
        f
        \colon
        x
        \longrightarrow
        y,
        \qquad
        g
        \colon
        y
        \longrightarrow
        z,
    \]
    there is a canonical equivalence
    \[
        \chi_{f,g}
        \colon
        (g\circ f)^R
        \xrightarrow{\;\simeq\;}
        f^R\circ g^R;
    \]

    \item the equivalences
    \[
        \chi_{f,g}
    \]
    are coherently associative and unital;

    \item right mating is compatible with vertical and horizontal
    composition of $2$-morphisms;

    \item the inverse equivalence is given by taking left adjoints and
    left mates.
\end{enumerate}
\end{theorem}

\begin{proof}
Let
\[
    *
\]
denote the terminal $(\infty,2)$-category. There are canonical
identifications
\[
    \Fun(*,\mathcal X)_{\mathrm{lax}}
    \simeq
    \mathcal X
    \simeq
    \Fun(*,\mathcal X)_{\mathrm{oplax}}.
\]
A functor
\[
    *
    \longrightarrow
    \mathcal X
\]
selects an object, a lax or oplax transformation selects a
$1$-morphism, and a modification selects a $2$-morphism.

The coherent mate theorem gives an equivalence
\[
    \Fun(*,\mathcal X)_{\mathrm{lax}}^{\mathrm{radj}}
    \simeq
    \left(
        \Fun(*,\mathcal X)_{\mathrm{oplax}}^{\mathrm{ladj}}
    \right)^{\mathrm{coop}}.
\]
Under the preceding identifications, this becomes
\[
    \mathcal X^{\mathrm R}
    \simeq
    \bigl(
        \mathcal X^{\mathrm L}
    \bigr)^{\mathrm{coop}}.
\]
Taking the inverse equivalence gives
\[
    R_{\mathcal X}
    \colon
    \bigl(
        \mathcal X^{\mathrm L}
    \bigr)^{\mathrm{coop}}
    \xrightarrow{\;\simeq\;}
    \mathcal X^{\mathrm R}.
\]
See
\cite[Theorem~5.3.6]{AbellanGagnaHaugsengMates}.

Since
\[
    R_{\mathcal X}
\]
is an $(\infty,2)$-functor, it preserves all composition operations in
its source. When its action is described using the underlying
$1$-morphisms and $2$-morphisms of
\[
    \mathcal X,
\]
the source coopposite reverses both levels.

Consequently, for composable left adjoints
\[
    f
    \colon
    x
    \longrightarrow
    y,
    \qquad
    g
    \colon
    y
    \longrightarrow
    z,
\]
the image of the composite is canonically equivalent to
\[
    f^R\circ g^R.
\]
This gives
\[
    \chi_{f,g}
    \colon
    (g\circ f)^R
    \xrightarrow{\;\simeq\;}
    f^R\circ g^R.
\]

The associativity pentagon and unit triangles for
\[
    \chi
\]
are the pseudofunctorial coherence diagrams of
\[
    R_{\mathcal X}.
\]
Compatibility of right mating with vertical and horizontal composition
follows from functoriality on $2$-morphisms and coherent pasting.

Taking left adjoints and left mates gives the inverse equivalence.
\end{proof}

\section{Specified adjunctions and the walking-adjunction resolution}
\label{sec:specified-adjunctions}

The coherent mate theorem makes adjoint formation functorial on the
sub-$(\infty,2)$-categories of adjointable morphisms. For the construction
of the complete weak involution, it is useful to retain a presentation in
which the right adjoints, units, and counits are part of the represented
diagram.

\begin{definition}[The walking adjunction]
\label{def:walking-adjunction}
Let
\[
    \mathsf{Adj}
\]
denote the walking adjunction. It has two distinguished objects
\[
    0,
    \qquad
    1,
\]
a distinguished $1$-morphism
\[
    \ell
    \colon
    0
    \longrightarrow
    1,
\]
a specified right adjoint
\[
    r
    \colon
    1
    \longrightarrow
    0,
\]
and coherent unit, counit, and triangle data
\[
    \eta
    \colon
    \id_0
    \Longrightarrow
    r\ell,
    \qquad
    \epsilon
    \colon
    \ell r
    \Longrightarrow
    \id_1.
\]

It is characterized by the universal property that, for every
$(\infty,2)$-category
\[
    \mathcal X,
\]
evaluation at
\[
    \ell
\]
identifies the space of functors
\[
    \mathsf{Adj}
    \longrightarrow
    \mathcal X
\]
with the space of $1$-morphisms of
\[
    \mathcal X
\]
equipped with complete coherent right-adjoint data.
\end{definition}

\begin{construction}[Walking strings of adjunctions]
\label{con:walking-strings-adjunctions}
For
\[
    n\geq 0,
\]
let
\[
    \mathsf{Adj}_{[n]}
\]
denote the walking string of
\[
    n
\]
composable adjunctions.

Thus
\[
    \mathsf{Adj}_{[0]}
\]
is the terminal $(\infty,2)$-category. Let
\[
    i_0,i_1
    \colon
    \mathsf{Adj}_{[0]}
    \longrightarrow
    \mathsf{Adj}
\]
denote the inclusions of the initial and terminal distinguished objects.
For
\[
    n\geq 1,
\]
there is an iterated pushout presentation
\[
\begin{aligned}
    \mathsf{Adj}_{[n]}
    \simeq{}&
    \mathsf{Adj}
    \amalg_{\mathsf{Adj}_{[0]}}
    \mathsf{Adj}
    \amalg_{\mathsf{Adj}_{[0]}}
    \cdots
    \amalg_{\mathsf{Adj}_{[0]}}
    \mathsf{Adj},
\end{aligned}
\]
with
\[
    n
\]
copies of
\[
    \mathsf{Adj}.
\]
At every pushout, the map from the preceding copy is
\[
    i_1,
\]
and the map to the following copy is
\[
    i_0.
\]
Thus the terminal distinguished object of each copy is identified with
the initial distinguished object of the next.

The coface maps compose adjacent specified adjunctions by pasting their
units and counits, and the codegeneracy maps insert identity adjunctions.
These maps make
\[
    [n]\longmapsto\mathsf{Adj}_{[n]}
\]
a cosimplicial $(\infty,2)$-category. Mapping this cosimplicial object
into an $(\infty,2)$-category produces the simplicial object which
encodes composition of specified adjunctions.
\end{construction}

\begin{theorem}[Specified-adjunction resolution]
\label{thm:specified-adjunction-resolution}
For every $(\infty,2)$-category
\[
    \mathcal X,
\]
there is an $(\infty,2)$-category
\[
    \operatorname{Adj}^{\mathrm L}(\mathcal X)
\]
of specified left adjunctions, natural in
\[
    \mathcal X,
\]
with the following properties.

\begin{enumerate}[label=(\roman*)]
    \item It has the same objects as
    \[
        \mathcal X.
    \]

    \item A $1$-morphism from
    \[
        x
    \]
    to
    \[
        y
    \]
    is complete specified adjunction data
    \[
        f
        \colon
        x
        \rightleftarrows
        y
        \colon
        g,
        \qquad
        f\dashv g.
    \]

    \item A $2$-morphism between specified adjunctions
    \[
        f\dashv g
        \qquad\text{and}\qquad
        f'\dashv g'
    \]
    is equivalently a $2$-morphism
    \[
        \alpha
        \colon
        f
        \Longrightarrow
        f'
    \]
    together with its coherently determined right mate
    \[
        \alpha^R
        \colon
        g'
        \Longrightarrow
        g.
    \]

    \item Composition is pasting of specified adjunctions. If
    \[
        f
        \colon
        x
        \rightleftarrows
        y
        \colon
        g
    \]
    and
    \[
        h
        \colon
        y
        \rightleftarrows
        z
        \colon
        k
    \]
    are specified adjunctions, their composite is
    \[
        hf
        \colon
        x
        \rightleftarrows
        z
        \colon
        gk.
    \]
    Its unit is
    \[
        \id_x
        \xRightarrow{\eta_f}
        gf
        \xRightarrow{g\eta_h f}
        gkhf,
    \]
    and its counit is
    \[
        hfgk
        \xRightarrow{h\epsilon_f k}
        hk
        \xRightarrow{\epsilon_h}
        \id_z.
    \]

    \item Forgetting the right adjoint, unit, counit, and mate data
    defines an equivalence
    \[
        p_{\mathrm L}
        \colon
        \operatorname{Adj}^{\mathrm L}(\mathcal X)
        \xrightarrow{\;\simeq\;}
        \mathcal X^{\mathrm L}.
    \]
\end{enumerate}
\end{theorem}

\begin{proof}
Choose a complete Segal-object presentation
\[
    N_\bullet\mathcal X
    \colon
    \Delta^{\op}
    \longrightarrow
    \Cat_\infty
\]
of
\[
    \mathcal X.
\]

Define
\[
    N_0^{\mathrm L}\mathcal X
    :=
    N_0\mathcal X,
\]
and let
\[
    N_1^{\mathrm L}\mathcal X
    \subseteq
    N_1\mathcal X
\]
be the full subcategory spanned by those $1$-morphisms of
\[
    \mathcal X
\]
which admit right adjoints.

For
\[
    n\geq 2,
\]
define
\[
\begin{aligned}
    N_n^{\mathrm L}\mathcal X
    :={}&
    N_1^{\mathrm L}\mathcal X
    \times_{N_0\mathcal X}
    \cdots
    \times_{N_0\mathcal X}
    N_1^{\mathrm L}\mathcal X,
\end{aligned}
\]
with
\[
    n
\]
factors.

Identity morphisms admit right adjoints, and left adjoints are closed
under composition. Hence the face and degeneracy maps of
\[
    N_\bullet\mathcal X
\]
restrict to
\[
    N_\bullet^{\mathrm L}\mathcal X.
\]
By construction this restricted simplicial object satisfies the Segal
condition.

Every equivalence in
\[
    \mathcal X
\]
is both a left and a right adjoint. Consequently,
\[
    N_\bullet^{\mathrm L}\mathcal X
\]
contains the complete equivalence subspace of
\[
    N_\bullet\mathcal X.
\]
It is therefore complete and presents the wide locally full
sub-$(\infty,2)$-category
\[
    \mathcal X^{\mathrm L}.
\]

We now construct a complete Segal object in which the adjunction data are
specified.

Set
\[
    \widetilde N_0^{\mathrm L}\mathcal X
    :=
    N_0\mathcal X.
\]

Let
\[
    \widetilde N_1^{\mathrm L}\mathcal X
\]
be the $\infty$-category whose objects are functors
\[
    H
    \colon
    \mathsf{Adj}
    \longrightarrow
    \mathcal X,
\]
and whose morphisms are colax transformations whose components at the
two distinguished objects of
\[
    \mathsf{Adj}
\]
are equivalences; higher morphisms are the corresponding modifications.

Evaluation at the distinguished left adjoint
\[
    \ell
    \colon
    0
    \longrightarrow
    1
\]
defines a functor
\[
    e_1
    \colon
    \widetilde N_1^{\mathrm L}\mathcal X
    \longrightarrow
    N_1^{\mathrm L}\mathcal X.
\]

We prove that
\[
    e_1
\]
is an equivalence of $\infty$-categories.

First, it is essentially surjective. A point of the fibre of
\[
    e_1
\]
over a left adjoint
\[
    f
    \colon
    x
    \longrightarrow
    y
\]
is a complete coherent choice of a right adjoint
\[
    g
    \colon
    y
    \longrightarrow
    x,
\]
together with unit, counit, triangle, and all higher adjunction data.
The walking-adjunction classification identifies this fibre with the
space of coherent adjunction extensions of
\[
    f.
\]
This space is nonempty and contractible. In a quasi-categorically
enriched presentation, this is
\cite[Theorem~4.4.18]{RiehlVerityHomotopyCoherentAdjunctions}.
Thus every left adjoint is in the essential image of
\[
    e_1.
\]

We next prove full faithfulness.

The icon identification
\cite[Theorem~7.3]{HaugsengLaxTransformations}
identifies the morphisms arising from the simplicial
$\infty$-categorical presentation of an $(\infty,2)$-category with
colax transformations whose components on objects are equivalences.

Consequently, for specified adjunctions
\[
    H
    =
    (f\dashv g)
    \qquad\text{and}\qquad
    H'
    =
    (f'\dashv g'),
\]
the mapping space
\[
    \Map_{\widetilde N_1^{\mathrm L}\mathcal X}(H,H')
\]
is the space of such colax transformations between the two
walking-adjunction diagrams.

Evaluation on the left legs gives a map
\[
\begin{aligned}
    \Map_{\widetilde N_1^{\mathrm L}\mathcal X}(H,H')
    \longrightarrow
    \Map_{N_1^{\mathrm L}\mathcal X}(f,f').
\end{aligned}
\]

The coherent mate theorem identifies a colax transformation between
specified adjunction diagrams with its left-leg component together with
the coherently determined right mate. More precisely, the
characterization of adjoints in lax and oplax functor
$(\infty,2)$-categories is supplied by
\cite[Corollary~5.2.10]{AbellanGagnaHaugsengMates},
and the equivalence on transformations and modifications is supplied by
\cite[Theorem~5.3.6]{AbellanGagnaHaugsengMates}.

It follows that the displayed map of mapping spaces is an equivalence.
Hence
\[
    e_1
    \colon
    \widetilde N_1^{\mathrm L}\mathcal X
    \xrightarrow{\;\simeq\;}
    N_1^{\mathrm L}\mathcal X
\]
is fully faithful and essentially surjective, and therefore an
equivalence of $\infty$-categories.

For
\[
    n\geq 2,
\]
define
\[
\begin{aligned}
    \widetilde N_n^{\mathrm L}\mathcal X
    :={}&
    \widetilde N_1^{\mathrm L}\mathcal X
    \times_{\widetilde N_0^{\mathrm L}\mathcal X}
    \cdots
    \times_{\widetilde N_0^{\mathrm L}\mathcal X}
    \widetilde N_1^{\mathrm L}\mathcal X,
\end{aligned}
\]
with
\[
    n
\]
factors.

By the pushout description of
\[
    \mathsf{Adj}_{[n]},
\]
this $\infty$-category is equivalently the category of functors
\[
    H
    \colon
    \mathsf{Adj}_{[n]}
    \longrightarrow
    \mathcal X
\]
and colax transformations whose components at all distinguished objects
are equivalences.

The coface and codegeneracy maps of
\Ref{con:walking-strings-adjunctions}
make
\[
    \widetilde N_\bullet^{\mathrm L}\mathcal X
\]
a simplicial object. Its inner coface maps paste adjacent specified
adjunctions, and its codegeneracy maps insert identity adjunctions.

Evaluation on the distinguished left legs defines a morphism of
simplicial objects
\[
    e_\bullet
    \colon
    \widetilde N_\bullet^{\mathrm L}\mathcal X
    \longrightarrow
    N_\bullet^{\mathrm L}\mathcal X.
\]

In degree zero,
\[
    e_0
\]
is the identity. In degree one,
\[
    e_1
\]
is the equivalence proved above. In every higher degree, both sides are
the corresponding iterated Segal fibre products of their degree-one
terms over their common degree-zero term. Hence
\[
    e_n
\]
is an equivalence for every
\[
    n\geq 0.
\]

Thus
\[
    e_\bullet
\]
is a degreewise equivalence of complete Segal objects.

Define
\[
    \operatorname{Adj}^{\mathrm L}(\mathcal X)
\]
to be the $(\infty,2)$-category presented by
\[
    \widetilde N_\bullet^{\mathrm L}\mathcal X.
\]
The degreewise equivalence induces
\[
    p_{\mathrm L}
    \colon
    \operatorname{Adj}^{\mathrm L}(\mathcal X)
    \xrightarrow{\;\simeq\;}
    \mathcal X^{\mathrm L}.
\]

The description of $1$-morphisms follows from the degree-one walking
adjunction.

The description of $2$-morphisms is the coherent mate equivalence
\[
    \alpha
    \colon
    f
    \Longrightarrow
    f'
    \qquad\longleftrightarrow\qquad
    \alpha^R
    \colon
    g'
    \Longrightarrow
    g.
\]

Compatibility with vertical and horizontal composition and with all
higher transformations is part of
\cite[Theorem~5.3.6]{AbellanGagnaHaugsengMates}.

Finally, the formulas for the unit and counit of a composite are the
standard pasting formulas and agree with the inner coface maps of the
walking-string construction.

Naturality in
\[
    \mathcal X
\]
is induced by postcomposition of walking-adjunction diagrams.
\end{proof}

\begin{remark}[Invariant use of the construction]
\label{rem:specified-adjunction-resolution-invariant}
The preceding proof uses a complete Segal-object presentation only to
construct the specified-adjunction resolution. Subsequent arguments use
only the invariant $(\infty,2)$-category
\[
    \operatorname{Adj}^{\mathrm L}(\mathcal X),
\]
the walking-adjunction description of its morphisms, and the equivalence
\[
    p_{\mathrm L}
    \colon
    \operatorname{Adj}^{\mathrm L}(\mathcal X)
    \xrightarrow{\;\simeq\;}
    \mathcal X^{\mathrm L}.
\]

In particular, the construction retains all equivalences between
objects: it does not replace the object $\infty$-groupoid of
\[
    \mathcal X
\]
by a discrete set.
\end{remark}

\section{The right-perfect Morita category and coherent adjoints}
\label{sec:right-perfect-morita-coherent-adjoints}

\begin{definition}[Right-perfect Morita category]
\label{def:right-perfect-morita-category}
Let
\[
    \Mor_A^{\mathrm{rperf}}
\]
denote the wide, locally full $(\infty,2)$-subcategory of
\[
    \Mor_A
\]
with:
\begin{enumerate}[label=(\roman*)]
    \item all associative $A$-algebras as objects;

    \item right-perfect bimodules as $1$-morphisms;

    \item all bimodule maps as $2$-morphisms.
\end{enumerate}
\end{definition}

\begin{proposition}[Right-perfect Morita morphisms and the mate calculus]
\label{prop:right-perfect-mate-calculus}
There is an equality of wide locally full subcategories
\[
    \Mor_A^{\mathrm{rperf}}
    =
    \Mor_A^{\mathrm L}.
\]
Consequently, coherent right adjoints and right mates on
\[
    \Mor_A^{\mathrm{rperf}}
\]
are supplied by
\[
    R_{\Mor_A}
    \colon
    \bigl(
        \Mor_A^{\mathrm{rperf}}
    \bigr)^{\mathrm{coop}}
    \xrightarrow{\;\simeq\;}
    \Mor_A^{\mathrm R}.
\]
\end{proposition}

\begin{proof}
By
\Ref{prop:right-perfect-morita-adjointability},
the right-perfect bimodules are precisely the $1$-morphisms admitting
Morita right adjoints. The result follows from
\Ref{thm:coherent-mate-calculus}.
\end{proof}

\begin{proposition}[Coherent adjoints of composites]
\label{prop:coherent-adjoints-composites}
For composable right-perfect bimodules
\[
    M
    \colon
    B
    \longrightarrow
    C,
    \qquad
    N
    \colon
    C
    \longrightarrow
    D,
\]
there is a canonical equivalence
\[
    \chi_{M,N}
    \colon
    (M\otimes_C N)^R
    \xrightarrow{\;\simeq\;}
    N^R\otimes_C M^R.
\]
The equivalences
\[
    \chi_{M,N}
\]
are coherently associative and unital.
\end{proposition}

\begin{proof}
The pointwise adjunction and its explicit unit and counit were constructed
in
\Ref{prop:morita-adjoint-composite}. Apply
\Ref{thm:coherent-mate-calculus}
to
\[
    \mathcal X=\Mor_A.
\]
The right adjoint of the Morita composite
\[
    M\otimes_C N
\]
is the reverse composite
\[
    N^R\otimes_C M^R.
\]
The associativity and unit coherences are the pseudofunctorial coherence
maps of the right-adjoint equivalence.
\end{proof}

\begin{proposition}[Closure of right-perfect Morita morphisms]
\label{prop:right-perfect-morita-closure}
The class of right-perfect bimodules contains the identity bimodules and
is closed under Morita composition.
\end{proposition}

\begin{proof}
The identity $B$--$B$-bimodule is right adjoint to itself and hence is
right perfect. Closure under composition follows from
\Ref{prop:morita-adjoint-composite}.
\end{proof}

\begin{proposition}[External tensor products of right-perfect bimodules]
\label{prop:external-tensor-right-perfect}
Let
\[
    M
    \colon
    B
    \longrightarrow
    C
\]
and
\[
    N
    \colon
    B'
    \longrightarrow
    C'
\]
be right perfect. Then
\[
    M\otimes_A N
    \colon
    B\otimes_A B'
    \longrightarrow
    C\otimes_A C'
\]
is right perfect, with right adjoint
\[
    M^R\otimes_A N^R.
\]
\end{proposition}

\begin{proof}
Tensoring the adjunctions
\[
    M\dashv M^R
    \qquad\text{and}\qquad
    N\dashv N^R
\]
and using the interchange equivalences of
\Ref{prop:symmetric-monoidal-higher-morita-category}
gives an adjunction
\[
    M\otimes_A N
    \dashv
    M^R\otimes_A N^R.
\]
Thus the external tensor product is right-adjointable and hence right
perfect.
\end{proof}

\section{Action reversal on Morita morphisms}
\label{sec:action-reversal-adjointable-morita}

\begin{proposition}[Coherent action-reversal involution]
\label{prop:action-reversal-equivalence}
Action reversal defines a coherent $1$-op involution
\[
    O_A
    \colon
    \Mor_A^{\op}
    \xrightarrow{\;\simeq\;}
    \Mor_A.
\]

On objects,
\[
    O_A(B)
    =
    B^{\op}.
\]

A $1$-morphism
\[
    B
    \longrightarrow
    C
\]
in
\[
    \Mor_A^{\op}
\]
is represented by a bimodule
\[
    M
    \colon
    C
    \longrightarrow
    B
\]
in
\[
    \Mor_A.
\]
Its image is
\[
    O_A(M)
    =
    M^{\op}
    \colon
    B^{\op}
    \longrightarrow
    C^{\op}.
\]

On $2$-morphisms,
\[
    O_A(\alpha)
    =
    \alpha^{\op}.
\]

The compositor is the canonical reversal equivalence
\[
    (M\otimes_C N)^{\op}
    \xrightarrow{\;\simeq\;}
    N^{\op}
    \otimes_{C^{\op}}
    M^{\op}.
\]

The involutive structure has underlying pseudonatural equivalence
\[
    \operatorname{bid}^{O_A}
    \colon
    \id_{\Mor_A}
    \xRightarrow{\;\simeq\;}
    O_A O_A^{\op},
\]
induced by the canonical double-opposite equivalences on algebras,
bimodules, and bimodule maps.

Together with its involutivity modification and higher coherences, these
data form a coherent homotopy-fixed-point structure for the
$1$-op action.

The equivalence restricts to
\[
    O_A
    \colon
    \bigl(
        \Mor_A^{\mathrm R}
    \bigr)^{\op}
    \xrightarrow{\;\simeq\;}
    \Mor_A^{\mathrm L}.
\]
\end{proposition}

\begin{proof}
Let
\[
    \sigma
    \colon
    \Assoc^\otimes
    \xrightarrow{\;\cong\;}
    \Assoc^\otimes
\]
be the canonical involution of the associative $\infty$-operad which
reverses every linear ordering. By
\cite[Remark~4.1.1.7]{LurieHigherAlgebra},
precomposition with
\[
    \sigma
\]
sends an associative algebra object to its opposite algebra and satisfies
\[
    \sigma^2
    =
    \id_{\Assoc^\otimes}.
\]

The coloured operads governing left modules, right modules, and
bimodules inherit compatible reflection involutions. On a
$B$--$C$-bimodule
\[
    {}_B M_C,
\]
the reflected structure is the
$C^{\op}$--$B^{\op}$-bimodule
\[
    {}_{C^{\op}}M^{\op}_{B^{\op}}
\]
whose underlying object is
\[
    M
\]
and whose actions are obtained by reversing the original right and left
actions.

Use the functorial complete Segal presentation of the higher Morita
category from
\cite[Remark~5.32 and Corollary~5.37]{HaugsengHigherMorita}.
At every simplicial degree, apply the reflection of the associative and
bimodule operads and reverse the order of the corresponding string of
composable bimodules.

On objects this sends
\[
    B
    \longmapsto
    B^{\op}.
\]
On a string
\[
    B_0
    \xrightarrow{M_1}
    B_1
    \xrightarrow{M_2}
    \cdots
    \xrightarrow{M_n}
    B_n,
\]
it gives
\[
    B_n^{\op}
    \xrightarrow{M_n^{\op}}
    B_{n-1}^{\op}
    \xrightarrow{M_{n-1}^{\op}}
    \cdots
    \xrightarrow{M_1^{\op}}
    B_0^{\op}.
\]

These degreewise operations commute with the face and degeneracy maps:
inner face maps are Morita composition, outer face maps forget an
endpoint, and degeneracy maps insert regular identity bimodules. The
reflection reverses inner composition and preserves identity
bimodules.

Consequently, before completion there is an order-two equivalence of
Segal objects
\[
    \widetilde O_A
    \colon
    \widetilde{\Mor}_A^{\op}
    \xrightarrow{\;\simeq\;}
    \widetilde{\Mor}_A.
\]
Because both the operadic reflection and reversal of strings square to
the identity, this equivalence is equipped at the presenting level with
a genuine coherent
\[
    C_2
\]
-action.

Passing to the functorial completion produces
\[
    O_A
    \colon
    \Mor_A^{\op}
    \xrightarrow{\;\simeq\;}
    \Mor_A
\]
together with the induced complete coherent
\[
    C_2
\]
-cocycle.

The compositor is the reversal equivalence constructed in
\Ref{prop:op-reversal-morita-composition}:
\[
    (M\otimes_C N)^{\op}
    \xrightarrow{\;\simeq\;}
    N^{\op}
    \otimes_{C^{\op}}
    M^{\op}.
\]
Its associativity and unit coherences are the coherent bar-reversal
compatibilities proved there.

Applying action reversal twice returns the original algebra, bimodule,
and bimodule map. At the presenting operadic level this is the identity
\[
    \sigma^2=\id,
\]
so after completion it supplies the pseudonatural equivalence
\[
    \operatorname{bid}^{O_A}
    \colon
    \id_{\Mor_A}
    \xRightarrow{\;\simeq\;}
    O_AO_A^{\op}
\]
together with its involutivity modification and all higher
\[
    C_2
\]
-coherences.

Finally, suppose
\[
    L\dashv M
\]
in
\[
    \Mor_A.
\]
Action reversal carries the unit and counit of this adjunction to the
unit and counit of
\[
    M^{\op}
    \dashv
    L^{\op}.
\]
Thus action reversal sends right adjoints to left adjoints and induces
the restricted equivalence
\[
    O_A
    \colon
    \bigl(
        \Mor_A^{\mathrm R}
    \bigr)^{\op}
    \xrightarrow{\;\simeq\;}
    \Mor_A^{\mathrm L}.
\]
\end{proof}

\begin{lemma}[Mate of the reversed mate]
\label{lem:mate-reversed-mate}
Let
\[
    \alpha
    \colon
    M
    \Longrightarrow
    N
\]
be a map between right-perfect bimodules. Relative to the reversed
adjunctions
\[
    (M^R)^{\op}
    \dashv
    M^{\op},
    \qquad
    (N^R)^{\op}
    \dashv
    N^{\op},
\]
the right mate of
\[
    (\alpha^R)^{\op}
    \colon
    (N^R)^{\op}
    \Longrightarrow
    (M^R)^{\op}
\]
is canonically equivalent to
\[
    \alpha^{\op}
    \colon
    M^{\op}
    \Longrightarrow
    N^{\op}.
\]

These equivalences are natural in
\[
    \alpha
\]
and coherently compatible with vertical composition, horizontal
composition, identity $2$-morphisms, and double action reversal.
\end{lemma}

\begin{proof}
The mate
\[
    \alpha^R
\]
is characterized by the counit equation
\[
    \epsilon_M^{C}
    \circ
    (\alpha^R\otimes_B\id_M)
    \simeq
    \epsilon_N^{C}
    \circ
    (\id_{N^R}\otimes_B\alpha).
\]
Applying the coherent action-reversal functor transforms this equation
into the counit characterization of the right mate of
\[
    (\alpha^R)^{\op}
\]
relative to the reversed adjunctions. The resulting mate is
\[
    \alpha^{\op}.
\]
Thus
\[
    \left(
        (\alpha^R)^{\op}
    \right)^R
    \simeq
    \alpha^{\op}.
\]

The compatibility with all compositions and higher pasting operations is
the compatibility of the coherent mate equivalence with postcomposition
by the $(\infty,2)$-functor
\[
    O_A.
\]
\end{proof}

\section{Invariant mate--reversal duality}
\label{sec:invariant-mate-reversal-duality}

The mate--reversal involution is first constructed on the
walking-adjunction resolution. There it is induced by an actual
involution of the indexing walking adjunction, so its complete
\[
    C_2
\]
-coherence is inherited rather than reconstructed from pointwise
biduality equivalences.

\begin{lemma}[Leg reversal of the walking adjunction]
\label{lem:walking-adjunction-leg-reversal}
There is a coherent involutive equivalence
\[
    s
    \colon
    \mathsf{Adj}
    \xrightarrow{\;\simeq\;}
    \mathsf{Adj}^{\op}
\]
which fixes the two distinguished objects and satisfies
\[
    s(\ell)=r^{\op},
    \qquad
    s(r)=\ell^{\op}.
\]
It carries the unit and counit of
\[
    \ell\dashv r
\]
to the unit and counit of
\[
    r^{\op}\dashv\ell^{\op}
\]
and, in the standard strict walking-adjunction model, satisfies
\[
    s^{\op}s
    =
    \id_{\mathsf{Adj}}.
\]
Invariantly, this equality supplies the canonical involutivity
equivalence and the complete
\[
    C_2
\]
-cocycle.

The construction extends levelwise to coherent involutive equivalences
\[
    s_{[n]}
    \colon
    \mathsf{Adj}_{[n]}
    \xrightarrow{\;\simeq\;}
    \mathsf{Adj}_{[n]}^{\op}
\]
compatible with the cosimplicial structure of
\Ref{con:walking-strings-adjunctions}.
\end{lemma}

\begin{proof}
Use the standard nerve-enriched walking adjunction model. By
\cite[Corollary~3.3.5]{RiehlVerityHomotopyCoherentAdjunctions}, it is
obtained by applying the nerve to the hom-categories of the
Schanuel--Street strict free adjunction $2$-category.

On the strict free adjunction $2$-category, define
\[
    s
    \colon
    \mathsf{Adj}
    \longrightarrow
    \mathsf{Adj}^{\op}
\]
by fixing the two objects and sending
\[
    \ell\longmapsto r^{\op},
    \qquad
    r\longmapsto\ell^{\op}.
\]
The original unit and counit have exactly the types of the unit and
counit of
\[
    r^{\op}\dashv\ell^{\op}
\]
in the $1$-opposite. The universal property of the strict free
adjunction therefore defines
\[
    s
\]
uniquely on the complete generating adjunction diagram.

Applying the same assignment twice fixes both objects, both generating
$1$-morphisms, the unit, and the counit. Hence
\[
    s^{\op}s
    =
    \id_{\mathsf{Adj}}
\]
strictly in this model. In particular, it supplies the complete coherent
\[
    C_2
\]
-cocycle after passage to the invariant $(\infty,2)$-category.

Apply the construction factorwise to the iterated pushout defining
\[
    \mathsf{Adj}_{[n]}.
\]
This gives strict involutive equivalences
\[
    s_{[n]}
    \colon
    \mathsf{Adj}_{[n]}
    \xrightarrow{\;\cong\;}
    \mathsf{Adj}_{[n]}^{\op}.
\]
The inner coface maps paste adjacent adjunctions, and the right adjoint of
a pasted composite is the reverse pasted composite. The outer cofaces
forget endpoints, and the codegeneracies insert identity adjunctions.
Therefore the maps
\[
    s_{[n]}
\]
commute with every coface and codegeneracy map.
\end{proof}

\begin{proposition}[Mate--reversal on specified adjunctions]
\label{prop:specified-mate-reversal-involution}
Let
\[
    \mathcal X
\]
be an $(\infty,2)$-category equipped with a coherent $1$-op involution
\[
    O
    \colon
    \mathcal X^{\op}
    \xrightarrow{\;\simeq\;}
    \mathcal X.
\]
Then there is a canonical coherent weak duality involution
\[
    J_O
    \colon
    \operatorname{Adj}^{\mathrm L}(\mathcal X)^{\mathrm{co}}
    \longrightarrow
    \operatorname{Adj}^{\mathrm L}(\mathcal X).
\]

For a specified adjunction represented by
\[
    H
    \colon
    \mathsf{Adj}
    \longrightarrow
    \mathcal X,
\]
it is represented by the composite
\[
\begin{aligned}
    \mathsf{Adj}
    &\xrightarrow{s}
    \mathsf{Adj}^{\op}
    \\
    &\xrightarrow{H^{\op}}
    \mathcal X^{\op}
    \\
    &\xrightarrow{O}
    \mathcal X.
\end{aligned}
\]
Thus, if
\[
    H
    =
    \bigl(
        f\dashv g
    \bigr),
\]
then
\[
    J_O(H)
    =
    \bigl(
        O(g)\dashv O(f)
    \bigr).
\]

For a $2$-morphism
\[
    \alpha
    \colon
    f
    \Longrightarrow
    f'
\]
with right mate
\[
    \alpha^R
    \colon
    g'
    \Longrightarrow
    g,
\]
the corresponding $2$-morphism in the co-opposite is sent to
\[
    O(\alpha^R)
    \colon
    O(g')
    \Longrightarrow
    O(g).
\]
\end{proposition}

\begin{proof}
By
\Ref{thm:specified-adjunction-resolution}, the specified-adjunction
category is presented by the complete Segal object obtained from the
walking strings
\[
    \mathsf{Adj}_{[n]}.
\]
In degree
\[
    n,
\]
define
\[
    J_O(H)
    :=
    O\circ H^{\op}\circ s_{[n]}.
\]
The compatibility of
\[
    s_{[n]}
\]
with cofaces and codegeneracies and the pseudofunctoriality of
\[
    O
\]
show that these assignments define an $(\infty,2)$-functor
\[
    J_O
    \colon
    \operatorname{Adj}^{\mathrm L}(\mathcal X)^{\mathrm{co}}
    \longrightarrow
    \operatorname{Adj}^{\mathrm L}(\mathcal X).
\]
On a single adjunction this is the displayed formula
\[
    f\dashv g
    \longmapsto
    O(g)\dashv O(f).
\]
On a $2$-morphism, the functor-category description and the coherent mate
theorem identify the left-leg component with
\[
    O(\alpha^R).
\]

Applying
\[
    J_O
\]
twice to a diagram
\[
    H
\]
gives
\[
\begin{aligned}
    J_OJ_O^{\mathrm{co}}(H)
    &\simeq
    O O^{\op}
    \circ
    H
    \circ
    s_{[n]}^{\op}s_{[n]}
    \\
    &\simeq
    O O^{\op}
    \circ
    H.
\end{aligned}
\]
The biduality of the coherent $1$-op involution,
\[
    \operatorname{bid}^{O}
    \colon
    \id_{\mathcal X}
    \xRightarrow{\;\simeq\;}
    O O^{\op},
\]
therefore induces by postcomposition a pseudonatural adjoint equivalence
\[
    \operatorname{bid}^{J}
    \colon
    \id_{\operatorname{Adj}^{\mathrm L}(\mathcal X)}
    \xRightarrow{\;\simeq\;}
    J_OJ_O^{\mathrm{co}}.
\]

The involutivity modification and all higher
\[
    C_2
\]
-cocycle simplices are obtained by applying the corresponding simplices
of the coherent action
\[
    O
\]
to the walking-adjunction diagrams and combining them with the strict
identity in the chosen walking-adjunction model
\[
    s_{[n]}^{\op}s_{[n]}
    =
    \id.
\]
Since both operations are maps of the complete Segal presentation, the
result is a homotopy-fixed-point structure for the co-action, not merely
a pointwise double-dual equivalence. Hence
\[
    J_O
\]
is a coherent weak duality involution.
\end{proof}

\begin{lemma}[Action reversal and adjointable morphisms]
\label{lem:one-op-reverses-adjunctions}
Let
\[
    \mathcal X
\]
be an $(\infty,2)$-category equipped with a coherent $1$-op involution
\[
    O
    \colon
    \mathcal X^{\op}
    \xrightarrow{\;\simeq\;}
    \mathcal X.
\]

If
\[
    f
    \colon
    x
    \longrightarrow
    y
\]
and
\[
    g
    \colon
    y
    \longrightarrow
    x
\]
form an adjunction
\[
    f\dashv g
\]
in
\[
    \mathcal X,
\]
then
\[
    O(g)\dashv O(f).
\]

Consequently,
\[
    O
\]
restricts to an equivalence
\[
    O^{\mathrm{adj}}
    \colon
    \bigl(
        \mathcal X^{\mathrm R}
    \bigr)^{\op}
    \xrightarrow{\;\simeq\;}
    \mathcal X^{\mathrm L}.
\]
\end{lemma}

\begin{proof}
An adjunction
\[
    f\dashv g
\]
consists of a unit and counit
\[
    \eta
    \colon
    \id_x
    \Longrightarrow
    gf,
    \qquad
    \epsilon
    \colon
    fg
    \Longrightarrow
    \id_y
\]
satisfying the triangle identities.

In
\[
    \mathcal X^{\op},
\]
the morphism
\[
    g
    \colon
    y
    \longrightarrow
    x
\]
is regarded as a morphism
\[
    x
    \longrightarrow
    y,
\]
and
\[
    f
\]
is regarded as a morphism
\[
    y
    \longrightarrow
    x.
\]
Since the $1$-opposite reverses $1$-morphisms but not $2$-morphisms, the
same unit and counit exhibit an adjunction
\[
    g\dashv f
\]
in
\[
    \mathcal X^{\op}.
\]

Applying the $(\infty,2)$-functor
\[
    O
    \colon
    \mathcal X^{\op}
    \longrightarrow
    \mathcal X
\]
gives an adjunction
\[
    O(g)\dashv O(f).
\]
The complete coherent unit, counit, and triangle data are transported by
functoriality.

Thus a right adjoint in
\[
    \mathcal X
\]
is carried, after passage to the $1$-opposite, to a left adjoint in
\[
    \mathcal X.
\]
This defines
\[
    O^{\mathrm{adj}}
    \colon
    \bigl(
        \mathcal X^{\mathrm R}
    \bigr)^{\op}
    \longrightarrow
    \mathcal X^{\mathrm L}.
\]

The inverse coherent $1$-op equivalence preserves adjunctions by the same
argument. Hence
\[
    O^{\mathrm{adj}}
\]
is an equivalence.
\end{proof}

\begin{theorem}[Abstract mate--reversal duality]
\label{thm:abstract-mate-reversal-duality}
Let
\[
    \mathcal X
\]
be an $(\infty,2)$-category equipped with a coherent $1$-op involution
\[
    O
    \colon
    \mathcal X^{\op}
    \xrightarrow{\;\simeq\;}
    \mathcal X.
\]

Let
\[
    R_{\mathcal X}
    \colon
    \bigl(
        \mathcal X^{\mathrm L}
    \bigr)^{\mathrm{coop}}
    \xrightarrow{\;\simeq\;}
    \mathcal X^{\mathrm R}
\]
be the coherent right-adjoint and right-mate equivalence of
\Ref{thm:coherent-mate-calculus}.

Taking the $1$-opposite gives
\[
    R_{\mathcal X}^{\op}
    \colon
    \bigl(
        \mathcal X^{\mathrm L}
    \bigr)^{\mathrm{co}}
    \xrightarrow{\;\simeq\;}
    \bigl(
        \mathcal X^{\mathrm R}
    \bigr)^{\op}.
\]
Using the restricted equivalence of
\Ref{lem:one-op-reverses-adjunctions},
define
\[
\begin{aligned}
    D_O
    &:=
    O^{\mathrm{adj}}
    \circ
    R_{\mathcal X}^{\op}
    \\
    &\colon
    \bigl(
        \mathcal X^{\mathrm L}
    \bigr)^{\mathrm{co}}
    \longrightarrow
    \mathcal X^{\mathrm L}.
\end{aligned}
\]

Then
\[
    D_O
\]
carries a canonical coherent weak duality involution.

On a left adjoint
\[
    f
    \colon
    x
    \longrightarrow
    y
\]
with right adjoint
\[
    f^R
    \colon
    y
    \longrightarrow
    x,
\]
one has
\[
    D_O(f)
    =
    O(f^R).
\]

For a $2$-morphism
\[
    \alpha
    \colon
    f
    \Longrightarrow
    g,
\]
the corresponding morphism in the co-opposite is
\[
    \alpha^{\mathrm{co}}
    \colon
    g
    \Longrightarrow
    f,
\]
and
\[
    D_O(\alpha^{\mathrm{co}})
    =
    O(\alpha^R)
    \colon
    O(g^R)
    \Longrightarrow
    O(f^R).
\]
\end{theorem}

\begin{proof}
By
\Ref{thm:specified-adjunction-resolution},
forgetting specified right adjoints, units, counits, and mate data gives
an equivalence
\[
    p_{\mathrm L}
    \colon
    \operatorname{Adj}^{\mathrm L}(\mathcal X)
    \xrightarrow{\;\simeq\;}
    \mathcal X^{\mathrm L}.
\]

By
\Ref{prop:specified-mate-reversal-involution},
the coherent $1$-op involution
\[
    O
\]
induces a coherent weak duality involution
\[
    J_O
    \colon
    \operatorname{Adj}^{\mathrm L}(\mathcal X)^{\mathrm{co}}
    \longrightarrow
    \operatorname{Adj}^{\mathrm L}(\mathcal X).
\]

Transport the complete homotopy-fixed-point structure of
\[
    J_O
\]
across the equivalence
\[
    p_{\mathrm L}.
\]
Transport of a homotopy-fixed-point structure along an equivalence is
canonical up to a contractible space of choices. It therefore produces
a coherent weak duality involution
\[
    D_O
    \colon
    \bigl(
        \mathcal X^{\mathrm L}
    \bigr)^{\mathrm{co}}
    \longrightarrow
    \mathcal X^{\mathrm L}.
\]

A specified adjunction
\[
    f\dashv f^R
\]
is sent by
\[
    J_O
\]
to
\[
    O(f^R)\dashv O(f).
\]
After forgetting the specified adjunction data, this gives
\[
    f
    \longmapsto
    O(f^R).
\]

Similarly, a $2$-morphism
\[
    \alpha
    \colon
    f
    \Longrightarrow
    g
\]
with right mate
\[
    \alpha^R
    \colon
    g^R
    \Longrightarrow
    f^R
\]
is sent, from the co-opposite, to
\[
    O(\alpha^R)
    \colon
    O(g^R)
    \Longrightarrow
    O(f^R).
\]

Under
\Ref{thm:coherent-mate-calculus},
these are precisely the objectwise and morphismwise formulas for
\[
    O^{\mathrm{adj}}
    \circ
    R_{\mathcal X}^{\op}.
\]
Thus the transported functor is the displayed composite
\[
    D_O.
\]

Its biduality pseudonatural equivalence, involutivity modification, and
all higher
\[
    C_2
\]
-cocycle data are transported from
\[
    J_O.
\]
Hence
\[
    D_O
\]
is a coherent weak duality involution.
\end{proof}

\section{Canonical Morita duality}
\label{sec:canonical-morita-duality}

Taking the $1$-opposite of the coherent right-adjoint equivalence gives
\[
    R_{\Mor_A}^{\op}
    \colon
    \bigl(
        \Mor_A^{\mathrm{rperf}}
    \bigr)^{\mathrm{co}}
    \xrightarrow{\;\simeq\;}
    \bigl(
        \Mor_A^{\mathrm R}
    \bigr)^{\op}.
\]

\begin{definition}[Canonical Morita duality]
\label{def:canonical-morita-duality}
Define
\[
    \mathbb D_A
    :=
    O_A\circ R_{\Mor_A}^{\op}.
\]
Thus
\[
    \mathbb D_A
    \colon
    \bigl(
        \Mor_A^{\mathrm{rperf}}
    \bigr)^{\mathrm{co}}
    \longrightarrow
    \Mor_A^{\mathrm{rperf}}.
\]

Explicitly,
\[
    \mathbb D_A(B)
    :=
    B^{\op},
\]
and for
\[
    M
    \colon
    B
    \longrightarrow
    C,
\]
\[
    \mathbb D_A(M)
    :=
    (M^R)^{\op}
    \colon
    B^{\op}
    \longrightarrow
    C^{\op}.
\]

For a $2$-morphism
\[
    \alpha
    \colon
    M
    \Longrightarrow
    N
\]
in
\[
    \Mor_A^{\mathrm{rperf}},
\]
the corresponding $2$-morphism
\[
    \alpha^{\mathrm{co}}
    \colon
    N
    \Longrightarrow
    M
\]
is sent to
\[
    \mathbb D_A(\alpha^{\mathrm{co}})
    :=
    (\alpha^R)^{\op}
    \colon
    (N^R)^{\op}
    \Longrightarrow
    (M^R)^{\op}.
\]
\end{definition}

\begin{theorem}[Coherent weak Morita duality involution]
\label{thm:weak-morita-duality-involution}
The canonical Morita duality
\[
    \mathbb D_A
    \colon
    \bigl(
        \Mor_A^{\mathrm{rperf}}
    \bigr)^{\mathrm{co}}
    \longrightarrow
    \Mor_A^{\mathrm{rperf}}
\]
is a coherent weak duality involution.

For composable right-perfect bimodules
\[
    M
    \colon
    B
    \longrightarrow
    C,
    \qquad
    N
    \colon
    C
    \longrightarrow
    D,
\]
there is a canonical compositor
\[
    \mu_{M,N}
    \colon
    \mathbb D_A(M\otimes_C N)
    \xrightarrow{\;\simeq\;}
    \mathbb D_A(M)
    \otimes_{C^{\op}}
    \mathbb D_A(N).
\]
These compositors are coherently associative and unital.

Define
\[
    \mathbb D_A^2
    :=
    \mathbb D_A
    \circ
    \mathbb D_A^{\mathrm{co}}.
\]
There is a pseudonatural adjoint equivalence
\[
    \operatorname{bid}^{\mathbb D}
    \colon
    \id_{\Mor_A^{\mathrm{rperf}}}
    \xRightarrow{\;\simeq\;}
    \mathbb D_A^2.
\]

For an associative $A$-algebra
\[
    B,
\]
its component is the regular Morita equivalence
\[
    \operatorname{bid}^{\mathbb D}_B
    \colon
    B
    \longrightarrow
    \mathbb D_A^2(B)
    \simeq
    B^{\op\op}
\]
associated to the canonical algebra equivalence
\[
    B
    \xrightarrow{\;\simeq\;}
    B^{\op\op}.
\]

For a right-perfect bimodule
\[
    M
    \colon
    B
    \longrightarrow
    C,
\]
pseudonaturality gives an invertible $2$-morphism
\[
\begin{aligned}
    \operatorname{bid}^{\mathbb D}_M
    \colon{}&
    \operatorname{bid}^{\mathbb D}_B
    \otimes_{\mathbb D_A^2(B)}
    \mathbb D_A^2(M)
    \\
    &\xRightarrow{\;\simeq\;}
    M
    \otimes_C
    \operatorname{bid}^{\mathbb D}_C.
\end{aligned}
\]
After transport along the endpoint equivalences, this comparison is the
canonical equivalence
\[
\begin{aligned}
    \mathbb D_A^2(M)
    &\simeq
    \left(
        \bigl((M^R)^{\op}\bigr)^R
    \right)^{\op}
    \\
    &\xrightarrow{\;\simeq\;}
    (M^{\op})^{\op}
    \\
    &\xrightarrow{\;\simeq\;}
    M.
\end{aligned}
\]

Finally, the biduality carries an invertible involutivity modification
comparing
\[
    \mathbb D_A
    \bigl(
        \operatorname{bid}^{\mathbb D}
    \bigr)
    \qquad\text{and}\qquad
    \operatorname{bid}^{\mathbb D}_{\mathbb D_A},
\]
and this modification belongs to the complete coherent
\[
    C_2
\]
-cocycle of the weak duality involution.
\end{theorem}

\begin{proof}
By
\Ref{prop:right-perfect-mate-calculus},
\[
    \Mor_A^{\mathrm{rperf}}
    =
    \Mor_A^{\mathrm L}.
\]
By
\Ref{prop:action-reversal-equivalence}, action reversal is a coherent
$1$-op involution
\[
    O_A
    \colon
    \Mor_A^{\op}
    \xrightarrow{\;\simeq\;}
    \Mor_A.
\]
Applying
\Ref{thm:abstract-mate-reversal-duality} with
\[
    \mathcal X
    =
    \Mor_A,
    \qquad
    O
    =
    O_A
\]
shows that
\[
    \mathbb D_A
    =
    O_A\circ R_{\Mor_A}^{\op}
\]
is a coherent weak duality involution, including its complete
biduality, modification, and higher
\[
    C_2
\]
-cocycle.

For composable right-perfect bimodules, the coherent right-adjoint
compositor gives
\[
    (M\otimes_C N)^R
    \xrightarrow{\;\simeq\;}
    N^R\otimes_C M^R.
\]
Applying action reversal and the compositor of
\[
    O_A
\]
gives
\[
\begin{aligned}
    \mathbb D_A(M\otimes_C N)
    &=
    \bigl(
        (M\otimes_C N)^R
    \bigr)^{\op}
    \\
    &\xrightarrow{\;\simeq\;}
    (N^R\otimes_C M^R)^{\op}
    \\
    &\xrightarrow{\;\simeq\;}
    (M^R)^{\op}
    \otimes_{C^{\op}}
    (N^R)^{\op}
    \\
    &=
    \mathbb D_A(M)
    \otimes_{C^{\op}}
    \mathbb D_A(N).
\end{aligned}
\]
This is
\[
    \mu_{M,N}.
\]
Its associativity and unit coherences are inherited from the
walking-adjunction pasting construction and the coherent action-reversal
compositors.

For the explicit biduality comparison, begin with a specified adjunction
\[
    M\dashv M^R.
\]
Action reversal gives the specified adjunction
\[
    (M^R)^{\op}
    \dashv
    M^{\op}.
\]
Thus, before forgetting the specified adjunction data, the selected right
adjoint of
\[
    (M^R)^{\op}
\]
is precisely
\[
    M^{\op}.
\]
After passing to any coherently selected right-adjoint functor, the
contractible uniqueness of right adjoints supplies a canonical
comparison
\[
    \rho_M
    \colon
    \left(
        (M^R)^{\op}
    \right)^R
    \xrightarrow{\;\simeq\;}
    M^{\op}.
\]
Action reversal then gives
\[
\begin{aligned}
    \mathbb D_A^2(M)
    &\simeq
    \left(
        \left(
            (M^R)^{\op}
        \right)^R
    \right)^{\op}
    \\
    &\xrightarrow{\rho_M^{\op}}
    (M^{\op})^{\op}
    \\
    &\xrightarrow{\;\simeq\;}
    M.
\end{aligned}
\]
For a bimodule map
\[
    \alpha
    \colon
    M
    \Longrightarrow
    N,
\]
\Ref{lem:mate-reversed-mate} gives
\[
    \left(
        \left(
            (\alpha^R)^{\op}
        \right)^R
    \right)^{\op}
    \simeq
    \alpha.
\]
These formulas are the visible components of the homotopy-fixed-point
structure already constructed on specified adjunctions. They are
therefore automatically compatible with all horizontal and vertical
pasting operations and with the complete
\[
    C_2
\]
-cocycle.
\end{proof}

\begin{proposition}[Symmetric monoidality of canonical Morita duality]
\label{prop:morita-duality-symmetric-monoidal}
The symmetric monoidal structure on
\[
    \Mor_A
\]
restricts to
\[
    \Mor_A^{\mathrm{rperf}}.
\]
With respect to this restricted structure, the canonical Morita duality
admits a canonical strong symmetric monoidal refinement
\[
    \mathbb D_A^\otimes
    \colon
    \bigl(
        \Mor_A^{\mathrm{rperf}}
    \bigr)^{\mathrm{co},\otimes}
    \longrightarrow
    \bigl(
        \Mor_A^{\mathrm{rperf}}
    \bigr)^\otimes.
\]

Its unit comparison is the canonical algebra equivalence
\[
    A^{\op}
    \xrightarrow{\;\simeq\;}
    A.
\]

For associative $A$-algebras
\[
    B
    \qquad\text{and}\qquad
    C,
\]
its objectwise tensor comparison is
\[
    \nu_{B,C}
    \colon
    (B\otimes_A C)^{\op}
    \xrightarrow{\;\simeq\;}
    B^{\op}\otimes_A C^{\op}.
\]

For right-perfect bimodules
\[
    M
    \colon
    B
    \longrightarrow
    C,
    \qquad
    N
    \colon
    B'
    \longrightarrow
    C',
\]
its tensor comparison is the canonical equivalence
\[
\begin{aligned}
    \nu_{M,N}
    \colon
    \mathbb D_A(M\otimes_A N)
    &=
    \bigl(
        (M\otimes_A N)^R
    \bigr)^{\op}
    \\
    &\xrightarrow{\;\simeq\;}
    (M^R\otimes_A N^R)^{\op}
    \\
    &\xrightarrow{\;\simeq\;}
    (M^R)^{\op}
    \otimes_A
    (N^R)^{\op}
    \\
    &=
    \mathbb D_A(M)\otimes_A\mathbb D_A(N).
\end{aligned}
\]

The comparisons
\[
    \nu
\]
are coherently associative, unital, and symmetric. They are compatible
with:
\begin{enumerate}[label=(\roman*)]
    \item the Morita compositors
    \[
        \mu_{M,N};
    \]

    \item the biduality pseudonatural equivalence
    \[
        \operatorname{bid}^{\mathbb D};
    \]

    \item the involutivity modification and complete coherent
    \[
        C_2
    \]
    -cocycle of
    \[
        \mathbb D_A.
    \]
\end{enumerate}
\end{proposition}

\begin{proof}
By
\Ref{prop:external-tensor-right-perfect},
the external tensor product of two right-perfect bimodules is right
perfect. Identity bimodules are right perfect. Hence
\[
    \Mor_A^{\mathrm{rperf}}
\]
is a symmetric monoidal wide locally full sub-$(\infty,2)$-category of
\[
    \Mor_A.
\]

Let
\[
    M\dashv M^R
    \qquad\text{and}\qquad
    N\dashv N^R
\]
be specified Morita adjunctions. Tensoring their units and counits and
using the interchange equivalences of
\Ref{prop:symmetric-monoidal-higher-morita-category}
gives a specified adjunction
\[
    M\otimes_A N
    \dashv
    M^R\otimes_A N^R.
\]
Explicitly, the unit is the composite
\[
\begin{aligned}
    B\otimes_A B'
    &\xrightarrow{
        \eta_M^B\otimes_A\eta_N^{B'}
    }
    (M\otimes_C M^R)
    \otimes_A
    (N\otimes_{C'}N^R)
    \\
    &\xrightarrow{\;\simeq\;}
    (M\otimes_A N)
    \otimes_{C\otimes_A C'}
    (M^R\otimes_A N^R),
\end{aligned}
\]
and the counit is obtained similarly from
\[
    \epsilon_M^C
    \qquad\text{and}\qquad
    \epsilon_N^{C'}.
\]
The triangle identities follow from the two original triangle identities
and the interchange coherences.

Consequently, the specified-adjunction resolution
\[
    \operatorname{Adj}^{\mathrm L}(\Mor_A)
\]
inherits a symmetric monoidal structure for which
\[
    p_{\mathrm L}
    \colon
    \operatorname{Adj}^{\mathrm L}(\Mor_A)
    \xrightarrow{\;\simeq\;}
    \Mor_A^{\mathrm{rperf}}
\]
is a symmetric monoidal equivalence.

The operadic reflection used to construct
\[
    O_A
\]
commutes with the external tensor product. Hence action reversal has a
canonical strong symmetric monoidal refinement. On algebras it gives
\[
    (B\otimes_A C)^{\op}
    \xrightarrow{\;\simeq\;}
    B^{\op}\otimes_A C^{\op},
\]
and on bimodules it gives
\[
    (P\otimes_A Q)^{\op}
    \xrightarrow{\;\simeq\;}
    P^{\op}\otimes_A Q^{\op}.
\]

The leg-reversal involution of the walking adjunction commutes with
tensoring specified adjunctions. Therefore the specified
mate--reversal involution
\[
    J_{O_A}
\]
is strong symmetric monoidal. Transporting this structure through
\[
    p_{\mathrm L}
\]
gives a strong symmetric monoidal structure on
\[
    \mathbb D_A.
\]

For the visible comparison on bimodules, the tensor product of specified
adjunctions gives the canonical right-adjoint comparison
\[
    (M\otimes_A N)^R
    \xrightarrow{\;\simeq\;}
    M^R\otimes_A N^R.
\]
Applying action reversal gives
\[
\begin{aligned}
    \bigl(
        (M\otimes_A N)^R
    \bigr)^{\op}
    &\xrightarrow{\;\simeq\;}
    (M^R\otimes_A N^R)^{\op}
    \\
    &\xrightarrow{\;\simeq\;}
    (M^R)^{\op}\otimes_A(N^R)^{\op}.
\end{aligned}
\]
This is
\[
    \nu_{M,N}.
\]

All associativity, unit, symmetry, and interchange diagrams hold on the
specified-adjunction resolution because they are obtained by tensoring
complete adjunction diagrams and applying the coherent operadic
reflection. The biduality, involutivity modification, and higher
\[
    C_2
\]
-cocycle are constructed on the same specified diagrams. The tensor
comparisons therefore intertwine the complete weak-duality structure,
not only its objectwise double-dual equivalences.
\end{proof}

\begin{remark}
\label{rem:op-versus-morita-duality}
The operation
\[
    M
    \longmapsto
    M^{\op}
\]
reverses the direction of $1$-morphisms.

The operation
\[
    M
    \longmapsto
    (M^R)^{\op}
\]
uses two reversals and therefore preserves the direction of
$1$-morphisms while reversing $2$-morphisms.

It is the latter operation that defines the coherent weak Morita duality
involution.
\end{remark}

\section{Base change of right-perfect Morita theory}
\label{sec:base-change-right-perfect-morita}

\begin{proposition}[Morita base change]
\label{prop:morita-base-change}
Let
\[
    f
    \colon
    A
    \longrightarrow
    A'
\]
be a morphism in
\[
    \CAlg(\V).
\]
Extension of scalars induces a strong symmetric monoidal
$(\infty,2)$-functor
\[
    F_f^{\mathrm{Mor}}
    \colon
    \Mor_A
    \longrightarrow
    \Mor_{A'}.
\]

On objects,
\[
    F_f^{\mathrm{Mor}}(B)
    =
    B\otimes_A A'.
\]
On bimodules,
\[
    F_f^{\mathrm{Mor}}(M)
    =
    M\otimes_A A'.
\]
On $2$-morphisms,
\[
    F_f^{\mathrm{Mor}}(\alpha)
    =
    \alpha\otimes_A A'.
\]

For composable bimodules
\[
    M
    \colon
    B
    \longrightarrow
    C,
    \qquad
    N
    \colon
    C
    \longrightarrow
    D,
\]
there is a canonical equivalence
\[
\begin{aligned}
    \phi_{f;M,N}
    \colon
    (M\otimes_C N)\otimes_A A'
    \xrightarrow{\;\simeq\;}
    (M\otimes_A A')
    \otimes_{C\otimes_A A'}
    (N\otimes_A A').
\end{aligned}
\]
These equivalences are coherently associative and unital.

For bimodules
\[
    M
    \colon
    B
    \longrightarrow
    C,
    \qquad
    N
    \colon
    B'
    \longrightarrow
    C',
\]
there is a canonical external tensor comparison
\[
\begin{aligned}
    \psi_{f;M,N}
    \colon
    (M\otimes_A N)\otimes_A A'
    \xrightarrow{\;\simeq\;}
    (M\otimes_A A')
    \otimes_{A'}
    (N\otimes_A A').
\end{aligned}
\]

If
\[
    A
    \xrightarrow{f}
    A'
    \xrightarrow{g}
    A''
\]
are composable, there is a pseudonatural adjoint equivalence
\[
    \gamma_{g,f}
    \colon
    F_g^{\mathrm{Mor}}
    \circ
    F_f^{\mathrm{Mor}}
    \xRightarrow{\;\simeq\;}
    F_{g\circ f}^{\mathrm{Mor}}
\]
whose component on every algebra or bimodule is induced by
\[
    (X\otimes_A A')\otimes_{A'}A''
    \xrightarrow{\;\simeq\;}
    X\otimes_A A''.
\]
There is also a unit equivalence
\[
    \gamma_A
    \colon
    F_{\id_A}^{\mathrm{Mor}}
    \xRightarrow{\;\simeq\;}
    \id_{\Mor_A}.
\]
The transformations
\[
    \gamma_{g,f}
    \qquad\text{and}\qquad
    \gamma_A
\]
satisfy the usual pseudofunctorial pentagon and triangle coherences.
\end{proposition}

\begin{proof}
Extension of scalars
\[
    -\otimes_A A'
    \colon
    \Mod_A(\V)
    \longrightarrow
    \Mod_{A'}(\V)
\]
is colimit-preserving and strong symmetric monoidal. It therefore sends
associative $A$-algebras to associative $A'$-algebras and sends
bimodules and bimodule maps to their base changes.

Let
\[
    M
    \colon
    B
    \longrightarrow
    C,
    \qquad
    N
    \colon
    C
    \longrightarrow
    D
\]
be composable. The relative tensor product is the geometric realization
\[
    M\otimes_C N
    \simeq
    \left|
        [n]
        \longmapsto
        M\otimes_A C^{\otimes_A n}\otimes_A N
    \right|.
\]
Since extension of scalars preserves geometric realizations,
\[
\begin{aligned}
    (M\otimes_C N)\otimes_A A'
    &\simeq
    \left|
        [n]
        \longmapsto
        \bigl(
            M\otimes_A C^{\otimes_A n}\otimes_A N
        \bigr)
        \otimes_A A'
    \right|.
\end{aligned}
\]
Strong symmetric monoidality gives degreewise equivalences
\[
\begin{aligned}
    &
    \bigl(
        M\otimes_A C^{\otimes_A n}\otimes_A N
    \bigr)
    \otimes_A A'
    \\
    &\qquad\simeq
    (M\otimes_A A')
    \otimes_{A'}
    (C\otimes_A A')^{\otimes_{A'}n}
    \otimes_{A'}
    (N\otimes_A A').
\end{aligned}
\]
The geometric realization of the simplicial object on the right is
\[
    (M\otimes_A A')
    \otimes_{C\otimes_A A'}
    (N\otimes_A A').
\]
This constructs
\[
    \phi_{f;M,N}.
\]

The comparisons for iterated Morita composition arise from the
corresponding degreewise comparisons of multisimplicial bar
constructions. Their associativity and unit coherences follow from the
associativity and unit coherences of extension of scalars together with
the canonical Fubini equivalences for geometric realizations.

The external tensor comparison
\[
    \psi_{f;M,N}
\]
is the strong symmetric monoidal comparison for
\[
    -\otimes_A A'.
\]
Its compatibility with Morita composition is the interchange coherence
between
\[
    \phi
    \qquad\text{and}\qquad
    \psi.
\]

Thus the assignments on algebras, bimodules, and bimodule maps define a
strong symmetric monoidal $(\infty,2)$-functor
\[
    F_f^{\mathrm{Mor}}.
\]

For composable base changes, associativity of relative tensor product
gives canonical equivalences
\[
    (X\otimes_A A')\otimes_{A'}A''
    \xrightarrow{\;\simeq\;}
    X\otimes_A A''
\]
natural for algebras, bimodules, and bimodule maps. Compatibility with
the bar comparisons makes these equivalences pseudonatural and
compatible with Morita composition. They therefore define
\[
    \gamma_{g,f}.
\]

The canonical equivalence
\[
    X\otimes_A A
    \simeq
    X
\]
defines
\[
    \gamma_A.
\]
The pentagon and triangle coherences are the associativity and unit
coherences of iterated extension of scalars. Hence base change is
pseudofunctorial in the commutative base.
\end{proof}

\begin{proposition}[Base change and coherent action reversal]
\label{prop:base-change-opposites}
Define
\[
    \mathcal O_f^{\mathrm L}
    :=
    O_{A'}
    \circ
    \bigl(F_f^{\mathrm{Mor}}\bigr)^{\op}
\]
and
\[
    \mathcal O_f^{\mathrm R}
    :=
    F_f^{\mathrm{Mor}}
    \circ
    O_A,
\]
viewed as functors
\[
    \Mor_A^{\op}
    \longrightarrow
    \Mor_{A'}.
\]
There is a pseudonatural adjoint equivalence
\[
    \omega_f
    \colon
    \mathcal O_f^{\mathrm L}
    \xRightarrow{\;\simeq\;}
    \mathcal O_f^{\mathrm R}.
\]

For every associative $A$-algebra
\[
    B,
\]
the object component is the regular Morita equivalence associated to the
canonical algebra equivalence
\[
    \omega_{f,B}^{\mathrm{alg}}
    \colon
    (B\otimes_A A')^{\op}
    \xrightarrow{\;\simeq\;}
    B^{\op}\otimes_A A'.
\]

Let
\[
    M
    \colon
    B
    \longrightarrow
    C
\]
be a bimodule, and regard it as a morphism
\[
    C
    \longrightarrow
    B
\]
in
\[
    \Mor_A^{\op}.
\]
The pseudonaturality constraint is an invertible $2$-morphism
\[
\begin{aligned}
    \Omega_{f,M}
    \colon{}&
    \omega_{f,C}
    \otimes_{\mathcal O_f^{\mathrm R}(C)}
    \mathcal O_f^{\mathrm R}(M)
    \\
    &\xRightarrow{\;\simeq\;}
    \mathcal O_f^{\mathrm L}(M)
    \otimes_{\mathcal O_f^{\mathrm L}(B)}
    \omega_{f,B}.
\end{aligned}
\]
After transporting the endpoint actions along
\[
    \omega_{f,B}^{\mathrm{alg}}
    \qquad\text{and}\qquad
    \omega_{f,C}^{\mathrm{alg}},
\]
this cell is represented by the canonical identification of the
underlying base-changed action-reversed objects
\[
    (M\otimes_A A')^{\op}
    \simeq
    M^{\op}\otimes_A A'.
\]

The pseudonatural equivalence
\[
    \omega_f
\]
is coherently compatible with relative tensor products, bimodule maps,
double action reversal, and the complete coherent
\[
    C_2
\]
-structures on
\[
    O_A
    \qquad\text{and}\qquad
    O_{A'}.
\]
\end{proposition}

\begin{proof}
Extension of scalars is a strong symmetric monoidal functor
\[
    -\otimes_A A'
    \colon
    \Mod_A(\V)
    \longrightarrow
    \Mod_{A'}(\V).
\]
It commutes with the order-reversal automorphism of the associative and
bimodule operads. Therefore functoriality of the higher Morita
construction produces a pseudonatural equivalence
\[
    O_{A'}
    \circ
    \bigl(F_f^{\mathrm{Mor}}\bigr)^{\op}
    \xRightarrow{\;\simeq\;}
    F_f^{\mathrm{Mor}}
    \circ
    O_A.
\]
On algebras this is induced by the identity on the underlying object and
the equality of the two reversed base-changed multiplication diagrams,
giving
\[
    \omega_{f,B}^{\mathrm{alg}}.
\]
The same comparison for the two module actions gives the displayed
pseudonaturality cell on bimodules.

At the bar-construction level, compatibility with composition is the
canonical equivalence
\[
    (M\otimes_C N)\otimes_A A'
    \xrightarrow{\;\simeq\;}
    (M\otimes_A A')
    \otimes_{C\otimes_A A'}
    (N\otimes_A A').
\]
The comparisons on all iterated bar constructions are induced by the
same symmetric monoidal interchange and order-reversal maps. They
therefore satisfy the pseudonaturality and compositor coherences.

Finally, since the comparison is induced before passing to algebras and
bimodules by a map of the order-two operadic actions, it intertwines the
double-opposite bidualities, involutivity modifications, and complete
\[
    C_2
\]
-cocycles.
\end{proof}

\begin{proposition}[Base change of Morita adjoints]
\label{prop:base-change-morita-adjoints}
Let
\[
    M
    \colon
    B
    \longrightarrow
    C
\]
be right perfect. Then
\[
    F_f^{\mathrm{Mor}}(M)
    =
    M\otimes_A A'
\]
is right perfect.

There is a canonical equivalence of right adjoints
\[
    \beta_{f,M}
    \colon
    (M\otimes_A A')^R
    \xrightarrow{\;\simeq\;}
    M^R\otimes_A A'.
\]

It is uniquely characterized by compatibility with the transported
adjunction
\[
    M\otimes_A A'
    \dashv
    M^R\otimes_A A'.
\]

These comparisons are coherently compatible with identity bimodules,
Morita composition, and right mates.
\end{proposition}

\begin{proof}
Applying
\[
    F_f^{\mathrm{Mor}}
\]
to the specified adjunction
\[
    M\dashv M^R
\]
gives an adjunction
\[
    M\otimes_A A'
    \dashv
    M^R\otimes_A A'.
\]
Thus the base-changed bimodule is right-adjointable and hence right
perfect.

The chosen right adjoint
\[
    (M\otimes_A A')^R
\]
and the transported right adjoint
\[
    M^R\otimes_A A'
\]
are right adjoints of the same $1$-morphism. The contractible uniqueness
of right adjoints gives a canonical equivalence
\[
    \beta_{f,M}
    \colon
    (M\otimes_A A')^R
    \xrightarrow{\;\simeq\;}
    M^R\otimes_A A'
\]
compatible with the units and counits.

For a bimodule map
\[
    \alpha
    \colon
    M
    \Longrightarrow
    N,
\]
compatibility with right mates is the canonical equivalence
\[
    F_f^{\mathrm{Mor}}(\alpha^R)
    \circ
    \beta_{f,N}
    \simeq
    \beta_{f,M}
    \circ
    \bigl(
        F_f^{\mathrm{Mor}}(\alpha)
    \bigr)^R.
\]

For composable right-perfect bimodules
\[
    M
    \colon
    B
    \longrightarrow
    C,
    \qquad
    N
    \colon
    C
    \longrightarrow
    D,
\]
the two composites comparing
\[
    \bigl(
        F_f^{\mathrm{Mor}}(M\otimes_C N)
    \bigr)^R
\]
with
\[
    F_f^{\mathrm{Mor}}(N^R)
    \otimes_{F_f^{\mathrm{Mor}}(C)}
    F_f^{\mathrm{Mor}}(M^R)
\]
are equivalences between right adjoints of the same composite and preserve
the same transported counit. They are therefore canonically equivalent.

This proves coherent compatibility with identities, composition, and
right mates.
\end{proof}

\begin{corollary}[Base-change functor on right-perfect Morita categories]
\label{cor:base-change-right-perfect-morita}
The functor
\[
    F_f^{\mathrm{Mor}}
\]
restricts to an $(\infty,2)$-functor
\[
    F_f^{\mathrm{rperf}}
    \colon
    \Mor_A^{\mathrm{rperf}}
    \longrightarrow
    \Mor_{A'}^{\mathrm{rperf}}.
\]
\end{corollary}

\begin{proof}
This follows from
\Ref{prop:base-change-morita-adjoints}.
\end{proof}

\begin{proposition}[Base change of canonical Morita duality]
\label{prop:base-change-morita-duality}
Define
\[
    L_f
    :=
    \mathbb D_{A'}
    \circ
    \bigl(F_f^{\mathrm{rperf}}\bigr)^{\mathrm{co}}
\]
and
\[
    R_f
    :=
    F_f^{\mathrm{rperf}}
    \circ
    \mathbb D_A.
\]
There is a pseudonatural adjoint equivalence
\[
    \kappa_f
    \colon
    L_f
    \xRightarrow{\;\simeq\;}
    R_f.
\]

For every associative $A$-algebra
\[
    B,
\]
the object component
\[
    \kappa_{f,B}
    \colon
    L_f(B)
    \longrightarrow
    R_f(B)
\]
is the regular Morita equivalence associated to
\[
    \omega_{f,B}^{\mathrm{alg}}
    \colon
    (B\otimes_A A')^{\op}
    \xrightarrow{\;\simeq\;}
    B^{\op}\otimes_A A'.
\]

For every right-perfect bimodule
\[
    M
    \colon
    B
    \longrightarrow
    C,
\]
the pseudonaturality constraint is an invertible $2$-morphism
\[
\begin{aligned}
    \kappa_{f,M}
    \colon{}&
    \kappa_{f,B}
    \otimes_{R_f(B)}
    R_f(M)
    \\
    &\xRightarrow{\;\simeq\;}
    L_f(M)
    \otimes_{L_f(C)}
    \kappa_{f,C}.
\end{aligned}
\]
It is the composite
\[
\begin{aligned}
    \kappa_{f,B}
    \otimes_{R_f(B)}
    F_f^{\mathrm{Mor}}
    \bigl(
        (M^R)^{\op}
    \bigr)
    &\xRightarrow{\;\Omega_{f,M^R}\;}
    \bigl(
        F_f^{\mathrm{Mor}}(M^R)
    \bigr)^{\op}
    \otimes_{L_f(C)}
    \kappa_{f,C}
    \\
    &\xRightarrow{
        \;(\beta_{f,M}^{\op})^{-1}
        \otimes
        \id
    \;}
    \bigl(
        F_f^{\mathrm{Mor}}(M)^R
    \bigr)^{\op}
    \otimes_{L_f(C)}
    \kappa_{f,C},
\end{aligned}
\]
where
\[
    \Omega_{f,M^R}
\]
is the pseudonaturality cell of
\[
    \omega_f
\]
from
\Ref{prop:base-change-opposites}, and
\[
    \beta_{f,M}
    \colon
    F_f^{\mathrm{Mor}}(M)^R
    \xrightarrow{\;\simeq\;}
    F_f^{\mathrm{Mor}}(M^R)
\]
is the comparison of right adjoints from
\Ref{prop:base-change-morita-adjoints}.

The pseudonatural equivalence
\[
    \kappa_f
\]
is compatible with the compositors, biduality pseudonatural
equivalences, involutivity modifications, and complete coherent
\[
    C_2
\]
-cocycles of
\[
    \mathbb D_A
    \qquad\text{and}\qquad
    \mathbb D_{A'}.
\]
Hence
\[
    F_f^{\mathrm{rperf}}
\]
is a morphism of $(\infty,2)$-categories with coherent weak duality
involution.
\end{proposition}

\begin{proof}
By naturality of the specified-adjunction resolution,
\[
    F_f^{\mathrm{Mor}}
\]
induces an $(\infty,2)$-functor
\[
    \widetilde F_f
    \colon
    \operatorname{Adj}^{\mathrm L}(\Mor_A)
    \longrightarrow
    \operatorname{Adj}^{\mathrm L}(\Mor_{A'}).
\]
It sends the complete specified adjunction
\[
    (M\dashv M^R,\eta_M,\epsilon_M)
\]
to
\[
    \bigl(
        F_f^{\mathrm{Mor}}(M)
        \dashv
        F_f^{\mathrm{Mor}}(M^R),
        F_f^{\mathrm{Mor}}(\eta_M),
        F_f^{\mathrm{Mor}}(\epsilon_M)
    \bigr).
\]

Let
\[
    J_A
\]
and
\[
    J_{A'}
\]
denote the specified mate--reversal involutions obtained from
\[
    O_A
    \qquad\text{and}\qquad
    O_{A'}
\]
by
\Ref{prop:specified-mate-reversal-involution}.
The pseudonatural equivalence
\[
    \omega_f
    \colon
    O_{A'}
    \circ
    \bigl(F_f^{\mathrm{Mor}}\bigr)^{\op}
    \xRightarrow{\;\simeq\;}
    F_f^{\mathrm{Mor}}
    \circ
    O_A
\]
acts on both legs, the unit, and the counit of every walking-adjunction
diagram. It therefore induces a pseudonatural adjoint equivalence
\[
    J_{A'}
    \circ
    \widetilde F_f^{\mathrm{co}}
    \xRightarrow{\;\simeq\;}
    \widetilde F_f
    \circ
    J_A
\]
on the specified-adjunction categories.

Transport this equivalence along the forgetful equivalences
\[
    \operatorname{Adj}^{\mathrm L}(\Mor_A)
    \xrightarrow{\;\simeq\;}
    \Mor_A^{\mathrm{rperf}}
\]
and
\[
    \operatorname{Adj}^{\mathrm L}(\Mor_{A'})
    \xrightarrow{\;\simeq\;}
    \Mor_{A'}^{\mathrm{rperf}}.
\]
The result is
\[
    \kappa_f
    \colon
    L_f
    \xRightarrow{\;\simeq\;}
    R_f.
\]
At an algebra
\[
    B,
\]
the component is the regular Morita equivalence associated to
\[
    \omega_{f,B}^{\mathrm{alg}}.
\]

Let
\[
    M
    \colon
    B
    \longrightarrow
    C
\]
be right perfect. Since
\[
    R_f(M)
    =
    F_f^{\mathrm{Mor}}
    \bigl(
        (M^R)^{\op}
    \bigr),
\]
the pseudonaturality cell
\[
    \Omega_{f,M^R}
\]
of
\[
    \omega_f
\]
gives
\[
\begin{aligned}
    \kappa_{f,B}
    \otimes_{R_f(B)}
    R_f(M)
    \xRightarrow{\;\simeq\;}
    \bigl(
        F_f^{\mathrm{Mor}}(M^R)
    \bigr)^{\op}
    \otimes_{L_f(C)}
    \kappa_{f,C}.
\end{aligned}
\]
The comparison of right adjoints
\[
    \beta_{f,M}
    \colon
    F_f^{\mathrm{Mor}}(M)^R
    \xrightarrow{\;\simeq\;}
    F_f^{\mathrm{Mor}}(M^R)
\]
after action reversal gives
\[
    \beta_{f,M}^{\op}
    \colon
    \bigl(
        F_f^{\mathrm{Mor}}(M)^R
    \bigr)^{\op}
    \xrightarrow{\;\simeq\;}
    \bigl(
        F_f^{\mathrm{Mor}}(M^R)
    \bigr)^{\op}.
\]
Postcomposing with its inverse gives exactly the displayed cell
\[
    \kappa_{f,M}.
\]
All endpoint algebras are now explicit, so this is a genuine Morita
$2$-morphism rather than an untyped equivalence between bimodules with
different actions.

Compatibility with $2$-morphisms is the mate identity
\[
    F_f^{\mathrm{Mor}}(\alpha^R)
    \circ
    \beta_{f,N}
    \simeq
    \beta_{f,M}
    \circ
    \bigl(
        F_f^{\mathrm{Mor}}(\alpha)
    \bigr)^R.
\]
Compatibility with horizontal composition follows from the coherent
right-adjoint compositors, the pseudonaturality and compositor coherence
of
\[
    \omega_f,
\]
and the interchange equivalence for extension of scalars.

Finally, the comparison was constructed on the complete
walking-adjunction diagrams. The pseudonatural equivalence
\[
    \omega_f
\]
intertwines the coherent $1$-op involutions before mate--reversal is
formed. Applying the functorial construction of
\Ref{prop:specified-mate-reversal-involution} therefore intertwines the
bidualities, involutivity modifications, and every higher simplex of the
\[
    C_2
\]
-cocycles. Transport along the forgetful equivalences preserves this
compatibility. Hence
\[
    F_f^{\mathrm{rperf}}
\]
is a morphism of coherent weak duality involutions.
\end{proof}

\begin{proposition}[Pseudofunctoriality of base change with Morita duality]
\label{prop:base-change-morita-duality-composition}
Let
\[
    A
    \xrightarrow{f}
    A'
    \xrightarrow{g}
    A''
\]
be composable morphisms in
\[
    \CAlg(\V).
\]
For brevity, write
\[
    F_f
    :=
    F_f^{\mathrm{rperf}},
    \qquad
    F_g
    :=
    F_g^{\mathrm{rperf}},
    \qquad
    F_{gf}
    :=
    F_{g\circ f}^{\mathrm{rperf}}.
\]

Let
\[
    \gamma_{g,f}
    \colon
    F_gF_f
    \xRightarrow{\;\simeq\;}
    F_{gf}
\]
be the base-change compositor induced by
\Ref{prop:morita-base-change}.

Define a pseudonatural adjoint equivalence
\[
\begin{aligned}
    \widetilde\kappa_{g,f}
    \colon
    \mathbb D_{A''}F_{gf}^{\mathrm{co}}
    \xRightarrow{\;\simeq\;}
    F_{gf}\mathbb D_A
\end{aligned}
\]
by the composite
\[
\begin{aligned}
    \mathbb D_{A''}F_{gf}^{\mathrm{co}}
    &\xRightarrow{
        \mathbb D_{A''}
        \ast
        \bigl(
            \gamma_{g,f}^{\mathrm{co}}
        \bigr)^{-1}
    }
    \mathbb D_{A''}F_g^{\mathrm{co}}F_f^{\mathrm{co}}
    \\
    &\xRightarrow{
        \kappa_g\ast F_f^{\mathrm{co}}
    }
    F_g\mathbb D_{A'}F_f^{\mathrm{co}}
    \\
    &\xRightarrow{
        F_g\ast\kappa_f
    }
    F_gF_f\mathbb D_A
    \\
    &\xRightarrow{
        \gamma_{g,f}\ast\mathbb D_A
    }
    F_{gf}\mathbb D_A.
\end{aligned}
\]

There is a canonical invertible modification
\[
    \Theta_{g,f}
    \colon
    \kappa_{g\circ f}
    \xRightarrow{\;\simeq\;}
    \widetilde\kappa_{g,f}.
\]

For the identity morphism
\[
    \id_A
    \colon
    A
    \longrightarrow
    A,
\]
the comparison
\[
    \kappa_{\id_A}
\]
identifies, under the unit equivalence
\[
    F_{\id_A}
    \xRightarrow{\;\simeq\;}
    \id_{\Mor_A^{\mathrm{rperf}}},
\]
with the identity pseudonatural transformation of
\[
    \mathbb D_A.
\]

For three composable morphisms
\[
    A
    \xrightarrow{f}
    A'
    \xrightarrow{g}
    A''
    \xrightarrow{h}
    A''',
\]
the two modifications obtained by applying
\[
    \Theta_{h,g},
    \qquad
    \Theta_{g,f},
    \qquad
    \Theta_{h,g\circ f},
    \qquad
    \Theta_{h\circ g,f}
\]
agree after transport along the associativity modification for the
base-change compositors
\[
    \gamma.
\]
The corresponding unit diagrams also commute.

Thus the assignment
\[
    A
    \longmapsto
    \bigl(
        \Mor_A^{\mathrm{rperf}},
        \mathbb D_A
    \bigr)
\]
together with
\[
    f
    \longmapsto
    \bigl(
        F_f^{\mathrm{rperf}},
        \kappa_f
    \bigr)
\]
is pseudofunctorial in the commutative base, with values in
$(\infty,2)$-categories equipped with coherent weak duality involution.
\end{proposition}

\begin{proof}
Work first on the specified-adjunction resolutions. Extension of scalars
induces
\[
    \widetilde F_f
    \colon
    \operatorname{Adj}^{\mathrm L}(\Mor_A)
    \longrightarrow
    \operatorname{Adj}^{\mathrm L}(\Mor_{A'}),
\]
and similarly for
\[
    g
    \qquad\text{and}\qquad
    g\circ f.
\]
By
\Ref{prop:morita-base-change},
there is a pseudonatural adjoint equivalence
\[
    \widetilde\gamma_{g,f}
    \colon
    \widetilde F_g\widetilde F_f
    \xRightarrow{\;\simeq\;}
    \widetilde F_{g\circ f}
\]
obtained by applying
\[
    (X\otimes_A A')\otimes_{A'}A''
    \xrightarrow{\;\simeq\;}
    X\otimes_A A''
\]
simultaneously to the left adjoint, right adjoint, unit, and counit of
every specified adjunction.

Let
\[
    J_A,
    \qquad
    J_{A'},
    \qquad
    J_{A''}
\]
be the specified mate--reversal involutions. The comparison
\[
    \omega_f
    \colon
    O_{A'}
    \circ
    \bigl(
        F_f^{\mathrm{Mor}}
    \bigr)^{\op}
    \xRightarrow{\;\simeq\;}
    F_f^{\mathrm{Mor}}
    \circ
    O_A
\]
induces
\[
    \widetilde\kappa_f
    \colon
    J_{A'}\widetilde F_f^{\mathrm{co}}
    \xRightarrow{\;\simeq\;}
    \widetilde F_fJ_A,
\]
and similarly for
\[
    g
    \qquad\text{and}\qquad
    g\circ f.
\]

All these comparisons are induced by extension of scalars and the same
order-reversal automorphism of the associative and bimodule operads.
For a specified walking-adjunction diagram
\[
    H
    \colon
    \mathsf{Adj}_{[n]}
    \longrightarrow
    \Mor_A,
\]
the two routes
\[
\begin{aligned}
    J_{A''}
    \widetilde F_{g\circ f}^{\mathrm{co}}(H)
    &\longrightarrow
    \widetilde F_{g\circ f}J_A(H)
\end{aligned}
\]
given respectively by
\[
    \widetilde\kappa_{g\circ f}
\]
and by the composite of
\[
    \widetilde\kappa_f,
    \qquad
    \widetilde\kappa_g,
    \qquad
    \widetilde\gamma_{g,f}
\]
are related by the associativity equivalence
\[
    (X\otimes_A A')\otimes_{A'}A''
    \xrightarrow{\;\simeq\;}
    X\otimes_A A''.
\]
This produces an invertible modification
\[
    \widetilde\Theta_{g,f}.
\]

Transport
\[
    \widetilde\Theta_{g,f}
\]
along the symmetric monoidal equivalences
\[
    \operatorname{Adj}^{\mathrm L}(\Mor_A)
    \xrightarrow{\;\simeq\;}
    \Mor_A^{\mathrm{rperf}},
\]
and the corresponding equivalences over
\[
    A'
    \qquad\text{and}\qquad
    A''.
\]
The resulting modification is
\[
    \Theta_{g,f}
    \colon
    \kappa_{g\circ f}
    \xRightarrow{\;\simeq\;}
    \widetilde\kappa_{g,f}.
\]

For the identity base change, the equivalence
\[
    X\otimes_A A
    \xrightarrow{\;\simeq\;}
    X
\]
acts identically on the specified left adjoint, right adjoint, unit, and
counit. Hence
\[
    \kappa_{\id_A}
\]
identifies with the identity transformation of
\[
    \mathbb D_A.
\]

For three composable base changes, both composites of the modifications
\[
    \widetilde\Theta
\]
are induced by the two standard reassociations of
\[
    \bigl(
        (X\otimes_A A')
        \otimes_{A'}A''
    \bigr)
    \otimes_{A''}A'''.
\]
They agree by the associativity pentagon for relative tensor product.
The unit coherences follow from its unit triangles.

All comparisons are constructed on complete walking-adjunction diagrams
and commute with action reversal there. They therefore intertwine the
biduality pseudonatural equivalences, involutivity modifications, and
every higher simplex of the coherent
\[
    C_2
\]
-cocycles. Transport along the specified-adjunction equivalences
preserves these coherences.

Hence base change together with
\[
    \kappa_f
\]
is pseudofunctorial as a morphism of coherent weak duality involutions.
\end{proof}

Canonical Morita duality changes every algebra endpoint to its opposite.
The next chapter equips selected endpoints with coherent equivalences to
their opposites. Those Hermitian fixed-point structures transport the
canonical Morita duality back to the original hom categories and produce
unitary bimodules and a global dagger.